\section{The local fundamental-group conjecture}\label{sec:F}
\subsection{The actual published boundary}
\begin{proposition}[F1: checked known cases; Known]\label{f:known}
Under the assumptions of Part F, the maximal prime-to-residue-characteristic
quotient is topologically finitely presented in the prime-to-characteristic
category. In residue characteristic zero this is the whole group.
In all residue characteristics the maximal pro-nilpotent quotient is
topologically finitely generated. In equal characteristic $p>0$, strong
$F$-regularity implies that the whole group is finite of order prime to $p$.
\end{proposition}
\begin{proof}[Verification of hypotheses and locators]
Grothendieck--Murre, \S9.1 and Theorem 9.2, p.~122, assume a complete
noetherian local ring of dimension two with algebraically closed residue
field, a closed subset $E$, and connected $S=\Spec(A)\setminus E$.
Their conclusion concerns the maximal prime-to-$p$ quotient; the
presentation convention on pp.~122--123 uses a free pro-prime-to-$p$
group. They allow $p=0$ \cite{GrothendieckMurre1971}.
The original two-page preview containing the theorem and definition was
opened directly. No analytic algebraization assumption is needed here.

Suzuki's Theorem 1.2, p.~2, proves the pro-nilpotent statement in equal
and mixed characteristic \cite{Suzuki2023}. Its hypotheses, printed as
Conjecture 1.1 on p.~1, require every irreducible component to have
dimension at least two and the punctured spectrum to be connected.
A normal two-dimensional local domain satisfies both conditions.
This theorem does not assert finite generation of the whole group.

Theorem A=Theorem 5.1 of \cite[p.~994]{CRST2018} assumes normality,
$F$-finiteness, strong $F$-regularity, strict henselianity, and dimension
at least two. It bounds the group order by $1/s(A)$ and makes it prime
to $p$. In the complete equicharacteristic situation here, strict
henselianity follows from the algebraically closed residue field.
A Cohen presentation makes $A$ a quotient of $k[[x_1,\ldots,x_e]]$;
Frobenius is finite on this power-series ring because $k$ is perfect,
so $A$ is $F$-finite. Strong $F$-regularity remains an additional
hypothesis and has not been deduced from normality.
\end{proof}

\begin{proposition}[F1b: all log terminal surfaces, including small and mixed characteristic; Known]
If the ring in Part F is log terminal, then its entire punctured
\'etale fundamental group is finite, in every equal or mixed
characteristic.
\end{proposition}
\begin{proof}[Verification of the exact published theorem]
This is Carvajal-Rojas--Yasuda, Theorem~B, p.~4 of the opened
27 March 2025 author version, equivalently Corollary~6.4,
pp.~29--30 \cite{CarvajalYasuda2025}. Its statement assumes exactly
a two-dimensional complete log terminal local ring with
algebraically closed residue field. The mixed-characteristic
setup on p.~28 specifies normality, a torsion canonical class
($\QQ$-Gorenstein), and positive log discrepancies on a minimal
resolution. The theorem includes characteristics $2$ and $3$ and
does not impose strong $F$-regularity. It asserts finiteness, not
that the order is prime to the residue characteristic.
\end{proof}

\begin{proposition}[F1c: the mixed-characteristic perfectoid-signature theorem; Known]
Let $A$ satisfy Part F in mixed characteristic $(0,p)$. If there is an
effective divisor $\Delta$ such that $K_A+\Delta$ is $\QQ$-Cartier
and $(A,\Delta)$ is BCM-regular, then the full punctured fundamental
group is finite, of order prime to $p$, and its order is at most the
inverse of the perfectoid signature of $A$.
\end{proposition}
\begin{proof}[Verification of the current primary source]
Cai, Lee, Ma, Schwede, and Tucker, Theorem~D on pp.~5--6 and
Theorem~7.0.2 on pp.~73--74 of the opened 3 July 2025 version,
give this assertion for any closed subset of codimension at least
two in a complete normal local domain of mixed characteristic with
algebraically closed residue field \cite{CaiEtAl2026}.
The journal article appeared online on 16 December 2025 and in
\emph{Inventiones mathematicae} 244 (2026), pp.~357--454.
Normality is not the BCM-regular-pair hypothesis. The latter is
used to make the perfectoid signature positive; the displayed
inverse-signature bound alone is vacuous when it is zero.
\end{proof}

The primary-source search through 20 November 2026 found no later
resolution for arbitrary normal complete surfaces. A negative search is not a
nonexistence proof. The full positive-residue-characteristic case
outside these additional regularity hypotheses remains
undecided in this work.

\begin{proposition}[F2: missing dimension qualification; Refuted wording]
\label{f:dimension}
The prompt's unqualified consequence for \emph{every} complete local
ring with algebraically closed residue field is false in dimension one.
\end{proposition}
\begin{proof}
Take $A=k[[t]]$ in characteristic $p>0$, with $k$ algebraically closed.
Artin--Schreier theory gives
\[
 H^1_{\mathrm{et}}(\Spec k((t)),\ZZ/p)
       =k((t))/(F-1)k((t)).
\]
The classes $t^{-n}$, $n>0$ and $p\nmid n$, are linearly independent
over $\mathbf F_p$. A nonzero finite combination has a largest pole
order prime to $p$, whereas a Laurent series $g$ with pole order $m$
has $g^p-g$ of pole order $pm$; an integral $g$ has no negative pole.
Thus this cohomology has infinite dimension. A group with $d$
topological generators has at most $p^d$ homomorphisms to $\ZZ/p$,
so its fundamental group is not finitely generated.
This does not refute the normal surface case. It explains why the
dimension-at-least-two qualification in SGA~2 XIII, Conjecture 3.1,
must remain in the consequence quoted from 58:19.
\end{proof}

\subsection{Why even prosolvable d\'evissage is insufficient}
\begin{proposition}[F3: an explicit metabelian obstruction; Proved]
\label{f:metabelian}
There is a metabelian profinite group $G$ whose maximal pro-nilpotent
quotient is $\ZZ_3$, such that every open subgroup has finitely generated
abelianization, but $G$ is not topologically finitely generated.
\end{proposition}
\begin{proof}
Let $f_n$ be the multiplicative order of $2$ modulo $3^n$ and
$L_n=\mathbf F_{2^{f_n}}$. Choose $\zeta_n\in L_n^\times$ of order
$3^n$. Put
\[
 V_n=L_n^n,\qquad V=\prod_{n\geq1}V_n,\qquad
 \Gamma=\ZZ_3,
\]
and let $\gamma\in\Gamma$ act on $V_n$ by multiplication by
$\zeta_n^{\gamma\bmod 3^n}$. This continuous action defines the
metabelian profinite group $G=V\rtimes\Gamma$.

The commutator operator with $1\in\Gamma$ is multiplication by
$\zeta_n-1$ on each coordinate and hence maps $V$ onto itself.
In any finite nilpotent quotient, the image of $V$ is a $2$-group
and that of $\Gamma$ a $3$-group. They commute, so the quotient kills
$[\Gamma,V]=V$. Conversely finite quotients of $\Gamma$ are
nilpotent. This proves the first assertion.

If $H$ is open, it contains $\prod_{n>M}V_n$ and $3^N\Gamma$ for
some $M,N$. Put $r=\max(M,N)$ and $T=\prod_{n>r}V_n$.
The commutator with $3^N\in\Gamma$ is invertible on every factor
of $T$, with continuous coordinatewise inverse. Thus
$T\subset[H,H]$. The quotient $H/T$ is a closed subgroup of
$(\prod_{n\leq r}V_n)\rtimes\Gamma$, with finite kernel and
procyclic image under projection to $\Gamma$. Finitely many
generators of the kernel and one lift of a generator of the image
generate it topologically. Its quotient $H^{\mathrm{ab}}$ is
therefore finitely generated. In particular
$H^1_{\mathrm{cont}}(H,\mathbf F_\ell)$ is finite for every prime
$\ell$.

Suppose $d$ elements generated $G$ densely. Their vector coordinates
in $V_n$ span an $L_n$-subspace $W_n$ of dimension at most $d$.
All words in their images in $V_n\rtimes(\Gamma/3^n\Gamma)$ lie
in $W_n\rtimes(\Gamma/3^n\Gamma)$, since the action is scalar.
For $n>d$ this is proper, contradicting generation of that finite
quotient. Thus $G$ is not finitely generated.
\end{proof}

This is a counterexample to an abstract inference, with no claimed
realization as a local fundamental group. It strengthens the
$\prod A_5$ warning in the supplied reading: the gap already occurs
for metabelian groups. The finite checks verify the relevant field
orders through $n=8$ and the perfection of $A_5$, while the preceding
proof handles the infinite product.

\paragraph{Approaches and exact gates.}
Rank-one Artin--Schreier/Frobenius estimates and their application to
each finite cover establish qualitative constant-coefficient
finiteness. Statement~\ref{f:metabelian} proves that these estimates
alone cannot give finite generation, even after restricting to
prosolvable quotients. Unbounded multiplicities of nontrivial finite
coefficient modules are an explicit remaining obstruction.
The resolution/intersection argument in Suzuki would need a uniform
extension to nonconstant local systems with controlled ramification;
no such estimate is proved here. Conversely wild covers of an affine
curve or cone do not automatically survive completion: regular
Artin--Schreier functions are killed by henselianity and the
algebraically closed residue field. No normal complete surface
counterexample has been constructed. The smallest undecided scope
recorded is a normal complete surface in positive residue
characteristic outside the log terminal and strongly $F$-regular
cases, with its full wild fundamental group.

\paragraph{The newer birational approach and its limitation.}
Carvajal-Rojas--Yasuda prove strict decrease of stringy motives
under ramified quasi-\'etale Galois covers and a descending chain
condition for the relevant log terminal surface invariants
\cite[Theorem 6.3 and Corollary 6.4, pp.~29--30]{CarvajalYasuda2025}.
This controls nonabelian wild covers in their stated class.
Positive discrepancies and the resulting convergence are
part of the construction. For an arbitrary normal non-log-terminal
surface, reusing that invariant and its descending chain condition
would be an additional unproved step; the theorem cannot be
applied after dropping log terminality.

\begin{corollary}[F3b: an explicit failure of uniform twisted-cohomology bounds; Proved]
For the group $G$ of Proposition~\ref{f:metabelian}, let $M_n=L_n$
with $V$ acting trivially and $\Gamma$ acting by $\zeta_n$.
Then
\[
 H^1_{\mathrm{cont}}(G,M_n)\simeq L_n^n,
 \qquad
 \frac{\dim_{\mathbf F_2}H^1(G,M_n)}{\dim_{\mathbf F_2}M_n}=n.
\]
By contrast, any profinite group generated by $d$ elements satisfies
$\dim_{\mathbf F_\ell}H^1(G,M)\leq d\dim_{\mathbf F_\ell}M$
for every finite continuous $\mathbf F_\ell[G]$-module $M$.
\end{corollary}
\begin{proof}
A continuous crossed homomorphism is determined by its values on a
topological generating set: the cocycle identity determines it on
words and continuity determines it on their closure. This proves
the stated necessary bound.

For our group, the restriction of a cocycle to $\Gamma$ is a
coboundary, since $\zeta_n-1$ is invertible on $L_n$: subtract
the coboundary whose value at $1\in\ZZ_3$ is the prescribed
value of the cocycle there, then use the density of the integers
in $\ZZ_3$. A normalized cocycle vanishing on $\Gamma$ is exactly
a continuous $\Gamma$-equivariant additive homomorphism $V\to L_n$.
Such a homomorphism factors through finitely many factors of $V$.
For $m<n$, the element $\gamma^{3^m}-1$ is zero on $L_m$
and invertible on $L_n$, forcing an equivariant homomorphism to
vanish. For $m>n$, use $\gamma^{3^n}-1$ in the reverse way. On $V_n=L_n^n$,
equivariance means $L_n$-linearity, because
$\mathbf F_2[\zeta_n]=L_n$. Thus the normalized cocycles form
$\Hom_{L_n}(L_n^n,L_n)=L_n^n$. No nonzero normalized cocycle
is a coboundary, since a coboundary vanishing on $\Gamma$ would
come from an element fixed by $\zeta_n$, necessarily zero.
This proves the formula and the unbounded ratio.
\end{proof}

\begin{proposition}[F4: a quantitative sufficient criterion for prosolvable quotients; Proved]
Let $G$ be a prosolvable profinite group. Suppose a real constant $C\geq0$
satisfies
\[
 \dim_{\mathbf F_\ell}H^1_{\mathrm{cont}}(G,M)
 \leq C\dim_{\mathbf F_\ell}M
\]
for every prime $\ell$ and every finite simple continuous
$\mathbf F_\ell[G]$-module $M$. Then $G$ has at most
$\lceil C\rceil+1$ topological generators. Conversely, $d$ generators
imply this bound with $C=d$.
\end{proposition}
\begin{proof}
Only sufficiency is new here. Put $D=\lceil C\rceil+1$.
For every finite quotient $F$ of $G$, inflation injects
$H^1(F,M)$ into $H^1(G,M)$ for an $F$-module $M$.
We show by induction on $|F|$ that every finite solvable quotient
$F$ is generated by $D$ elements, allowing identities in the tuple.
Choose a minimal nontrivial normal subgroup $M$ of $F$.
Solvability and minimality imply that $M$ is an elementary abelian
$\ell$-group and a simple module for $H=F/M$ under conjugation:
its derived subgroup, primary parts, and subgroup $\ell M$ are
characteristic, so minimality successively forces these assertions.
By induction choose a generating $D$-tuple of $H$.

Any lift of this tuple generates a subgroup $J$ with $JM=F$.
The intersection $J\cap M$ is an $H$-submodule, hence is either
$0$ or $M$. In the latter case $J=F$; in the former case $J$ is
a complement to $M$. Thus a nonsplit extension has a generating
lift immediately. In the split case identify $F=M\rtimes H$.
The complements are precisely the graphs of the crossed
homomorphisms $H\to M$. For the chosen tuple each complement
contains exactly one lifting tuple, so the number of nongenerating
lifts is $|Z^1(H,M)|$.

Restriction to $M$ and the splitting give, directly on cocycles,
\[
 H^1(F,M)\simeq H^1(H,M)\oplus\operatorname{End}_H(M).
\]
Indeed a cocycle vanishing on $H$ is exactly an $H$-equivariant
additive map $M\to M$, and a coboundary restricts to zero on $M$.
The endomorphism group contains its nonzero identity, so has
cardinality at least $\ell$. Consequently
\[
 |Z^1(H,M)|
 =|H^1(H,M)|\frac{|M|}{|M^H|}
 \leq \frac{|M|^{C+1}}{\ell}
 <|M|^D.
\]
There are $|M|^D$ lifting tuples, so a generating one exists.
This proves the finite induction.

For each open normal subgroup $N$ of $G$, the tuples in $G^D$
generating $G/N$ form a nonempty closed subset. Any finite
intersection contains the subset corresponding to the
intersection of the relevant normal subgroups, which is nonempty
by the finite result. Compactness therefore supplies one tuple
generating every finite quotient; its closed generated subgroup
is $G$. Necessity was proved in F3b by evaluating cocycles on a
generating tuple.
\end{proof}

\paragraph{The resulting geometric target.}
For the maximal prosolvable quotient of the local fundamental group,
it would suffice to prove a constant $C_A$, independent of the
finite simple local system $\mathcal M$ and its rank, with
\[
 \dim_{\mathbf F_\ell}H^1_{\mathrm{et}}
       (\Spec(A)\setminus\{\mathfrak m\},\mathcal M)
 \leq C_A\operatorname{rank}(\mathcal M).
\]
One only needs local systems whose monodromy is solvable; inflation
then bounds the cohomology of the prosolvable quotient. The argument
does not provide this uniform estimate. In particular, separate
finiteness on each finite cover has no uniform constant, as F3
demonstrates. Nonabelian simple chief factors of the full fundamental
group require an additional argument even after this target is met.

\subsection{A uniform bound for abelian prime-to-$p$ monodromy}

\begin{proposition}[F5: uniform twisted cohomology in a geometric subclass; Proved]
Let $A$ be a complete normal two-dimensional local domain of
characteristic $p>0$ with algebraically closed residue field $k$,
and put $U=\Spec A\setminus\{\mathfrak m\}$. There is a constant
$C_A$ such that, for every finite locally constant
$\mathbf F_p$-sheaf $M$ on $U$ whose monodromy group is abelian
of order prime to $p$,
\[
 \dim_{\mathbf F_p}H^1_{\mathrm{et}}(U,M)
 \leq C_A\operatorname{rank}_{\mathbf F_p}M.
\]
Consequently the group in F3 cannot be a quotient of such a local
fundamental group when $p=2$.
\end{proposition}

\begin{proof}
We first produce one finite cover, independent of $M$, on whose
fixed resolution all these coefficient systems extend without
ramification. We then make Suzuki's equal-characteristic
nearby-cycle argument uniform for the resulting torsion line bundles.
The argument uses the explicit filtration and Frobenius-cokernel
calculation of \cite[Propositions 2.1 and 3.1--3.4, pp.~4--7]{Suzuki2023};
the uniformity and the preliminary fixed-cover reduction are proved below.

\paragraph{A finite bound for abelian tame inertia.}
Choose a coefficient field $k\hookrightarrow A$ by the equal-characteristic
Cohen structure theorem. Since $k$ is perfect, a power-series presentation
makes $A$, and every scheme of finite type over it used below, $F$-finite.
Choose a resolution $f:X\to\Spec A$ with reduced exceptional divisor
$D=\bigcup_{i=1}^s D_i$ a simple normal crossing divisor, each $D_i$
smooth and $D$ nonempty; if $A$ is regular, first blow up its closed point.
Its intersection matrix $B=(D_i\cdot D_j)$ is negative definite,
so $b=|\det B|$ is a positive integer. Write
$Q=\pi_1(U)^{\mathrm{ab},(p')}$ and let $I_i\subset Q$ be the image
of tame inertia at $D_i$. These are procyclic groups.

For every character $\chi:Q\to\mu_n$ with $(n,p)=1$, Kummer theory
represents the torsor by a line bundle $L$ on $U$ with a trivialization
of $L^{\otimes n}$. Extend $L$ to a line bundle $\widetilde L$ on the
regular scheme $X$: extend a representing divisor and use regularity.
The trivialization identifies
\[
 \widetilde L^{\otimes n}\simeq
 \mathcal O_X\Bigl(\sum_i a_iD_i\Bigr)
\]
for integers $a_i$. Its inertia character at $D_i$ is $a_i\bmod n$,
with the usual common choice of roots of unity. Restriction to $D_j$
gives
\[
 \sum_i a_i(D_i\cdot D_j)
       =n\deg(\widetilde L|_{D_j}).
\]
The adjugate matrix therefore gives $ba_i=0\bmod n$ for every $i$.
Characters separate points of the abelian pro-prime-to-$p$ group $Q$,
so $bI_i=0$. The closed subgroup $I=\sum_i I_i$ is finite.

The same subgroup contains the image of tame inertia for every
exceptional divisorial valuation over $\Spec A$. Indeed, at the center
of that valuation on $X$, trivialize $\widetilde L$. The Kummer equation
has the form
\[
 z^n=u\prod_{i\in S}t_i^{a_i},
\]
where $u$ is a unit and the $t_i$ are the local equations of the
components through the center. Its valuation is
$\sum_{i\in S}v(t_i)a_i$. The unit contributes only unramified
monodromy. In particular, if $\chi$ is trivial on $I$, then all $a_i$ are zero
modulo $n$, so this valuation inertia also has trivial character.
Characters of $Q/I$ separate points; this places its image in $I$
without making compatible choices of inertia generators.
There is no wild inertia in an abelian prime-to-$p$ quotient.

Choose an open subgroup $Q'\subset Q$ with $Q'\cap I=0$; a finite
quotient separating the finitely many nonzero elements of $I$ does
this. Let $U'\to U$ be the connected finite abelian prime-to-$p$
cover defined by $Q'$. Normalize $\Spec A$ in its function field to
obtain $\Spec A'$; excellence makes this normalization finite, and
completeness and the algebraically closed residue field make $A'$
a complete normal local domain with residue field $k$.
Choose an SNC resolution $X'\to\Spec A'$ with reduced exceptional
curve $D'$. For every $M$ in the proposition, its pullback to $U'$
has trivial inertia at every exceptional divisor of $X'$: its image
in $Q$ lies both in $I$ and in $Q'$. Here a divisorial valuation of
the finite field extension restricts to an exceptional divisorial
valuation of the base field, because residue transcendence degree one
is preserved under finite extension. Normalize $X'$ in its finite
monodromy cover. It is unramified in codimension one, hence is
\'etale everywhere by purity of the branch locus on the regular
scheme $X'$ \cite[Lemma 58.21.4, tag 0BMB]{StacksBranchPurity}.
Thus every such $M|_{U'}$ extends to a locally constant
sheaf on the same fixed $X'$.

Restriction
$H^1(U,M)\to H^1(U',M|_{U'})$ is injective, since restriction followed
by corestriction is multiplication by the degree, which is prime to
$p$. It suffices to obtain a uniform bound on $U'$. In the rest of
the proof rename $A',X',D',U'$ as $A,X,D,U$; these are now fixed.

\paragraph{The extended coefficient bundle.}
Let $M$ extend to $X$, set $r=\operatorname{rank}_{\mathbf F_p}M$, and
let $V=M\otimes_{\mathbf F_p}\mathcal O_X$. This is a vector bundle
with its natural Frobenius-semilinear map $F_V$, whose linearization
$F_X^*V\to V$ is an isomorphism. Because the monodromy is finite,
abelian, and of order prime to $p$, the representation over $k$
splits into characters. Hence
\[
 V\simeq\bigoplus_{a=1}^r L_a,
 \qquad L_a\text{ a finite-order line bundle on }X.
\]
In particular $\deg(L_a|_{D_i})=0$ for every $a,i$.
The summands need not be individually Frobenius-stable.

Choose a positive exceptional cycle $Z=\sum_i c_iD_i$ such that
$J=\mathcal O_X(-Z)$ has positive degree on every $D_i$, as in
\cite[p.~6]{Suzuki2023}. All choices below depend on this fixed
resolution and $J$, not on $M$.
With $j:U\hookrightarrow X$, $i:D\hookrightarrow X$, and
$\Psi=i^*j_*$ on \'etale sheaves, the Artin--Schreier sequence for
$V$ gives
\[
 0\longrightarrow V(*D)/V
 \xrightarrow{F_V-1} V(*D)/V
 \longrightarrow R^1\Psi M\longrightarrow0.
\]
Here we identify coherent quotients supported on $D$ with their
pullbacks by $i$. Exactness follows locally from the ordinary
Artin--Schreier sequence after an \'etale trivialization of $M$;
$j$ is affine and $F_V-1$ is surjective on the extended bundle in
the \'etale topology.

Let $U_mR^1\Psi M$ be the image of $J^{-pm}V/V$.
Valuations along the regular components of $D$, and the fact that
$F_X^*V\to V$ is an isomorphism, give the exact sequence
\[
 0\longrightarrow J^{-m}V/V
 \xrightarrow{F_V-1}J^{-pm}V/V
 \longrightarrow U_mR^1\Psi M\longrightarrow0.
 \tag{F5.1}
\]
In detail, a section with a pole of order $e>0$ has Frobenius pole
order $pe$, so its leading pole cannot cancel with its own order-$e$
term; regularity tests membership in these divisor lattices at
codimension-one points. This also proves the claimed kernel and
preimage assertions in (F5.1).

\paragraph{A pole cutoff independent of the coefficients.}
For $m\geq2$, the same quotient calculation as in Suzuki's
Proposition~3.1, using the projection formula for Frobenius and
$F_X^*V\simeq V$, gives
\[
 \operatorname{gr}_m R^1\Psi M
 \simeq B_Z\otimes_{\mathcal O_Z}V|_Z
                  \otimes_{\mathcal O_Z}J^{-m}|_Z,
 \qquad
 B_Z=(F_*\mathcal O_{pZ})/\mathcal O_Z.
 \tag{F5.2}
\]
The sheaf $B_Z$ is locally free on $Z$: the finite flat Frobenius
on the regular scheme $X$ makes $(F_*\mathcal O_X)/\mathcal O_X$
locally free, and the displayed sheaf is its base change to $Z$.

Fix a filtration of $\mathcal O_Z$ with successive terms line bundles
$P_\nu$ on components $D_{i(\nu)}$, obtained by removing one
component at a time from the effective cycle $Z$. Tensoring by the
locally free sheaves in (F5.2) gives successive terms
\[
 (B_Z|_{D_{i(\nu)}}\otimes P_\nu)
   \otimes L_a|_{D_{i(\nu)}}
   \otimes J^{-m}|_{D_{i(\nu)}}.
\]
For each of the finitely many fixed vector bundles
$B_Z|_{D_i}\otimes P_\nu$, choose an injection into a finite sum
of copies of a line bundle $H_\nu$ of positive degree. Such an
injection is obtained by globally generating a sufficiently ample
twist of the dual bundle and then dualizing. Since
$\deg L_a|_{D_i}=0$ and $\deg J|_{D_i}>0$, all the resulting line
bundles have negative degree for $m\geq m_0$, with one $m_0$
independent of $a,M,r$. Therefore
\[
 H^0(D,\operatorname{gr}_mR^1\Psi M)=0
 \quad(m\geq m_0).
\]
Filtered colimits and quasi-compactness now give
\[
 H^0(D,R^1\Psi M)=H^0(D,U_{m_0}R^1\Psi M),
 \tag{F5.3}
\]
increasing $m_0$ by one if necessary.

\paragraph{The remaining pole range has dimension at most a constant times rank.}
Suzuki's negative-degree filtration, Proposition~2.1(1), applies to
$J^{-pm}/\mathcal O_X$ and $J^{-pm}/J^{-m}$. Tensoring every layer
by a finite-order $L_a$ keeps its degree negative. Thus both
quotients tensored by $V$ have no global sections. Taking cohomology
of (F5.1) injects $H^0(D,U_mR^1\Psi M)$ into the kernel of
\[
 F_V-1:H^1(D,J^{-m}V/V)\longrightarrow
                    H^1(D,J^{-pm}V/V),
\]
and the map denoted $1$ here is injective. If a Frobenius-linear
map and an injective linear map have common domain of $k$-dimension
$N$, their difference has an $\mathbf F_p$-kernel of dimension at
most $N$: project to a linear splitting of the latter map and use
the fixed-space bound for a semilinear endomorphism. For completeness,
$\mathbf F_p$-linearly independent fixed vectors are $k$-linearly
independent: normalize a shortest $k$-linear relation to have first
coefficient $1$, apply Frobenius, and subtract; minimality forces
all coefficients into $\mathbf F_p$, a contradiction. This is the
quantitative form of Suzuki's Lemma~3.3.

A fixed filtration of $J^{-m_0}/\mathcal O_X$ has line-bundle layers
$T_\nu$ on smooth components $D_{i(\nu)}$. Their degrees are
negative. On a smooth projective curve $C$ of genus $g$, Riemann--Roch
therefore gives, for every degree-zero $L$,
\[
 h^1(C,T_\nu\otimes L)=g-1-\deg T_\nu.
\]
The finite filtration and the splitting of $V$ imply
\[
 \dim_k H^1(D,J^{-m_0}V/V)\leq C_1r
\]
for a constant $C_1$ fixed by $X,J,m_0$. Combining this with
(F5.3) gives
\[
 \dim_{\mathbf F_p}H^0(D,R^1\Psi M)\leq C_1r.
\]

Finally, the five-term nearby-cycle sequence and proper base change
bound $H^1(U,M)$ by $H^1(D,M)$ and $H^0(D,R^1\Psi M)$.
The Artin--Schreier sequence on the proper curve $D$ identifies
$H^1(D,M)$ with the Frobenius-fixed subgroup of $H^1(D,V|_D)$: on the finite-dimensional space of
global sections, $F_V-1$ is surjective over the algebraically closed
field. The normalization sequence of the nodal curve $D$ bounds
\[
 h^1(D,L_a|_D)\leq\sum_i g(D_i)+\#\{\text{nodes of }D\},
\]
since each restriction to a component has degree zero. Consequently
$\dim_{\mathbf F_p}H^1(D,M)\leq C_2r$. This yields the desired
bound with $C_A=C_1+C_2$ for the fixed cover.

For the last assertion, suppose $\pi_1(U)$ surjected onto the group
of F3 in characteristic $2$. Its modules $M_n$ have cyclic odd-order
monodromy. Inflation and F3b would give
\[
 n=\frac{\dim H^1(G,M_n)}{\dim M_n}
 \leq\frac{\dim H^1(U,M_n)}{\dim M_n}\leq C_A
\]
for every $n$, a contradiction.
\end{proof}

\paragraph{Exact scope of this advance.}
F5 bounds coefficients with abelian prime-to-$p$ monodromy and
excludes the specific metabelian obstruction F3 as a geometric
quotient. It does not bound all simple modules of a prosolvable
fundamental group: their finite monodromy can be nonabelian and
can have order divisible by $p$. It therefore does not yet prove
finite generation of the maximal prosolvable quotient, much less
the full local fundamental group. The fixed-cover argument depends
on the finite abelian inertia subgroup and does not extend formally
to those other monodromy groups.

\begin{corollary}[F6: nonabelian pro-$p$ kernels over abelian tame quotients; Proved]
Under the equal-characteristic hypotheses of F5, every quotient $G$
of $\pi_1(U)$ which admits an exact sequence
\[
 1\longrightarrow P\longrightarrow G\longrightarrow\Gamma
 \longrightarrow1,
 \qquad P\text{ pro-}p,\quad
 \Gamma\text{ abelian pro-prime-to-}p,
\]
is topologically finitely generated. In fact one bound on the number
of generators works for all such quotients of the fixed $\pi_1(U)$.
The kernel $P$ is allowed to be nonabelian. The same finite-generation
conclusion holds when $\Gamma$ is a virtually abelian pro-prime-to-$p$
group, without asserting a uniform bound independent of its abelian
subgroup's index.
\end{corollary}
\begin{proof}
Let $d_0$ generate $\pi_1(U)^{\mathrm{ab},(p')}$, which is finitely
generated by F1. Then every $\Gamma$ in the first assertion has at
most $d_0$ generators. We bound $H^1(G,M)$ for every finite simple
$\mathbf F_\ell[G]$-module $M$.

For $\ell=p$, the normal pro-$p$ subgroup $P$ acts trivially on $M$.
Indeed its action factors through a finite $p$-group, whose every
nonzero characteristic-$p$ representation has a nonzero invariant
vector (induct on the group order using a central element of order
$p$ and the nonzero kernel of its nilpotent operator $z-1$).
The invariant subspace $M^P$ is $G$-stable, hence is all of the
simple module $M$. Thus $M$ has abelian prime-to-$p$ monodromy.
Inflation from $G$ to $\pi_1(U)$ is injective in degree one, and F5
therefore gives
\[
 \dim_{\mathbf F_p}H^1(G,M)\leq C_A\dim_{\mathbf F_p}M.
\]

For $\ell\ne p$, one has $H^1(P,M)=0$. A continuous cocycle factors
through a finite $p$-group together with its finite coefficient
action, and averaging by that group's invertible order makes it a
coboundary. Inflation--restriction now gives
\[
 H^1(G,M)\simeq H^1(\Gamma,M^P).
\]
Again $M^P$ is either $0$ or $M$ by simplicity and normality. In the
latter case evaluation on $d_0$ generators bounds the dimension of
this cohomology by $d_0\dim M$; in the former it is zero.

All finite quotients of $G$ are solvable, since they are extensions
of finite $p$-groups by finite abelian groups. Apply F4 with
$C=\max(C_A,d_0)$ to obtain at most $\lceil C\rceil+1$ generators,
uniformly in $G$.

If $\Gamma$ has an open abelian subgroup $\Gamma'$, its inverse
image $G'$ has finite index in $G$. The corresponding finite \'etale
cover $U'\to U$ has degree prime to $p$ and is the punctured spectrum
of a complete normal local surface $A'$ with the same algebraically
closed residue field. The first assertion applied to $A'$ shows
that $G'$ is finitely generated. Adding finitely many coset
representatives generates $G$.
\end{proof}

The class in F6 does not contain arbitrary solvable quotients:
alternating wild and tame normal factors can leave no normal
pro-$p$ subgroup with abelian tame quotient. Thus the full
prosolvable and full profinite problems remain separate gates.

\begin{proposition}[F7: arbitrary coefficients extending over a fixed resolution; Proved]
Let $A,U$ satisfy F5, and fix an SNC resolution $X\to\Spec A$,
with nonempty reduced exceptional curve $D$. There is a constant
$C_X$ such that every finite locally constant $\mathbf F_p$-sheaf
$M$ on $X$ satisfies
\[
 \dim_{\mathbf F_p}H^1(U,M|_U)
       \leq C_X\operatorname{rank}_{\mathbf F_p}M.
\]
The monodromy of $M$ may be nonabelian and may have order divisible
by $p$.
\end{proposition}
\begin{proof}
We give all changes needed to the uniform part of F5, retaining its
notation $V=M\otimes_{\mathbf F_p}\mathcal O_X$, $J$, $Z$, $B_Z$,
and the nearby-cycle filtration. The Artin--Schreier and Frobenius
identities (F5.1)--(F5.2) remain valid: they only use that $M$ is
locally constant on $X$, equivalently that $V$ has its unit
Frobenius structure. We must replace the decomposition into torsion
line bundles by rank-uniform estimates for $V$ itself.

For each smooth exceptional component $C=D_i$, the bundle $V|_C$
is trivialized by a finite \'etale cover of $C$. It is semistable of
degree zero. Indeed its determinant becomes trivial on that cover,
so has degree zero. The pullback of any subbundle becomes a
subbundle of a trivial bundle. Taking its determinant gives a
nonzero map to a finite sum of trivial line bundles, so its degree
is nonpositive. Degree multiplies by the positive covering degree,
so the original subbundle also has nonpositive degree. This proof
works equally for covers of degree divisible by $p$.

It follows that
\[
 H^0(C,V|_C\otimes T)=0\quad\text{if }\deg T<0,
 \qquad h^0(C,V|_C)\leq r,
 \quad r=\operatorname{rank}V.
\]
For the second inequality, evaluation of sections at any closed
point of $C$ is injective: a section vanishing there would give a
line subbundle of positive degree after saturation. Riemann--Roch
therefore gives
\[
 h^1(C,V|_C)\leq r g(C),\qquad
 h^1(C,V|_C\otimes T)=r(g(C)-1-\deg T)\quad(\deg T<0).
 \tag{F7.1}
\]

For the high-pole calculation in F5, replace a layer indexed by
$L_a$ by
\[
 V|_C\otimes(B_Z|_C\otimes P_\nu)\otimes J^{-m}|_C.
\]
The fixed vector bundle in parentheses embeds in a finite sum of
copies of the fixed line bundle $H_\nu$. Semistability and the first
vanishing above show that the displayed layer has no sections as
soon as $\deg H_\nu-m\deg J|_C<0$. Thus the same fixed $m_0$
works, independently of $V$ and $r$.

For the low-pole calculation, the negative-degree filtration of
$J^{-m}/\mathcal O_X$ and $J^{-pm}/J^{-m}$, tensored with $V$,
has no sections. Formula (F7.1) bounds its first cohomology by a
constant times $r$. The semilinear fixed-space lemma in F5 bounds
$H^0(D,R^1\Psi M)$ by the same type of constant. Finally the
normalization sequence for $D$ and (F7.1) give
\[
 h^1(D,V|_D)\leq r\left(\sum_i g(D_i)+\#\{\text{nodes}\}\right).
\]
As in F5, the Artin--Schreier sequence identifies $H^1(D,M)$ with
the Frobenius-fixed subgroup of $H^1(D,V|_D)$, so its
$\mathbf F_p$-dimension is at most the latter $k$-dimension.
The nearby-cycle five-term sequence finishes the proof.
\end{proof}

\begin{lemma}[F7a: a relative pro-$p$ generation criterion; Proved]
Let $1\to P\to G\to\Gamma\to1$ be an exact sequence of profinite
groups, with $P$ pro-$p$ and $\Gamma$ generated by $d$ elements.
Suppose that for every finite simple $\mathbf F_p[G]$-module $M$
there is a uniform bound
$\dim H^1(G,M)\leq C\dim M$. Then $G$ is generated by at most
$\max(d,\lceil C\rceil+1)$ elements. No solvability hypothesis on
$\Gamma$ is needed.
\end{lemma}
\begin{proof}
Put $D_0=\max(d,\lceil C\rceil+1)$. In any finite quotient $F$
of $G$, let $N$ be the image of $P$. Induct on $|N|$ to generate
$F$ by $D_0$ elements. If $N=1$, $F$ is a quotient of $\Gamma$
and the assertion holds. Otherwise choose a minimal nontrivial
$F$-normal subgroup $M\subseteq N$. It is an elementary abelian
$p$-group and a simple $\mathbf F_p[F/M]$-module, by the same
characteristic-subgroup argument as in F4. Induction generates
$F/M$ by $D_0$ elements.

The nonsplit case and the count of nongenerating lifts as
$|Z^1(F/M,M)|$ are exactly those in F4. In the split case the
nonzero $\operatorname{End}_{F/M}(M)$ summand in $H^1(F,M)$
gives the strict bound
\[
 |Z^1(F/M,M)|\leq |M|^{C+1}/p<|M|^{D_0}.
\]
Inflation to $G$ supplies the assumed bound on $H^1(F,M)$.
Thus a generating lift exists. Compactness of $G^{D_0}$, applied
to all finite quotients, completes the argument as in F4.
\end{proof}

\begin{corollary}[F7b: pro-$p$ extensions of quotients from the resolution; Proved]
Let $G$ be a quotient of $\pi_1(U)$ with normal pro-$p$ subgroup
$P$. Suppose the composite
\[
 \pi_1(U)\longrightarrow G\longrightarrow\Gamma=G/P
\]
factors through $\pi_1(X)$. Then $G$ is topologically finitely
generated, with a bound depending only on the fixed resolution $X$.
The group $\Gamma$ need not be abelian, solvable, or prime to $p$.
\end{corollary}
\begin{proof}
Every finite simple $\mathbf F_p[G]$-module is trivial on the normal
pro-$p$ subgroup, by the invariant-vector argument in F6. Its
pullback to $\pi_1(U)$ therefore extends to $X$ under the stated
factorization. Inflation and F7 give the uniform bound required
by F7a.

It remains to bound the generators of $\Gamma$. Proper
specialization over the complete henselian ring $A$ gives
$\pi_1(X)\simeq\pi_1(D)$; passing to the reduced special fibre does
not change finite \'etale covers
\cite[Lemma 58.9.1, tag 0A48]{StacksProperEtaleSpecialization}.
The connected curve $D$ is proper over the algebraically closed
field, so its fundamental group is finitely generated. For a
current stronger primary statement, see
\cite[Theorem 1.2, p.~1]{LaraSrinivasStix2024}, which proves finite
presentation for every connected proper scheme over an
algebraically closed field (and recalls SGA~1 X, Theorem~2.9 for
finite generation). Hence one fixed generator bound works for
all quotients $\Gamma$ of $\pi_1(X)$, and F7a applies.
\end{proof}

This does not assert that the kernel of $\pi_1(U)\to\pi_1(X)$ is
pro-$p$. Tame inertia and more complicated wild quotients may
remain in that kernel. The explicit normal-pro-$p$ hypothesis in
F7b is essential to the argument.

\subsection{Uniform tame coefficients: a root-stack argument}

\begin{proposition}[F8: uniform tame coefficient bound; Proved]
Let $A$ be a complete normal two-dimensional local domain of
characteristic $p>0$ with algebraically closed residue field $k$.
Put $U=\Spec A\setminus\{\mathfrak m\}$, and fix an SNC resolution
$X\to\Spec A$ with reduced exceptional divisor
$D=\bigcup_{i=1}^s C_i$. There is a constant $C_X$ such that
\[
 \dim_{\mathbf F_p}H^1(U,M)\leq C_X\operatorname{rank}M
 \tag{F8.0}
\]
for every finite locally constant $\mathbf F_p$-sheaf $M$ on $U$
whose monodromy is tamely ramified along $D$.
The finite monodromy group need not be abelian or of order prime to $p$.
\end{proposition}

\paragraph{Status.}
The complete proof passed an independent root audit, including the
primary sources and nearby-cycle formulas; BCE separately checked
the root-gerbe degree, coarse-descent, and character-count steps.
This is a proved result of the present work, without a historical
priority claim. The uniformity in root orders is part of the proof.

\begin{proof}
Choose a coefficient field $k\hookrightarrow A$. Cohen structure
expresses $A$ as a quotient of a power-series ring over perfect $k$,
which is finite over its Frobenius image; hence $A$ and the
essentially finite-type chart rings below are $F$-finite. If $A$ is regular, first blow up its closed point,
so $D$ is nonempty. Write
\[
 B_{ij}=C_i\cdot C_j,\qquad g_i=g(C_i),\qquad
 \nu_i=\#\{\text{crossings on }C_i\}.
\]
Choose positive integers $c_i$ so that
$J=\mathcal O_X(-Z)$, $Z=\sum_i c_iC_i$, has positive degree
\[
 d_i=-\sum_jc_j B_{ij}>0
\]
on every component. Such a cycle exists by the negative definiteness
used in Suzuki's proof. Set
\[
 L_i=\sum_jc_j|B_{ij}|,\qquad
 K_i=\max\{0,-B_{ii},2g_i-2+\nu_i\}.
 \tag{F8.1}
\]
These constants depend only on the chosen resolution.

\paragraph{1. Root stacks, extension, and degree conventions.}
Choose an integer $N$ prime to $p$ divisible by the tame inertia
orders of the finite monodromy of $M$ along all $C_i$. Let
\[
 \pi:\mathfrak X_N=\sqrt[N]{(X,C_1)}\mathbin{\times_X}\cdots
                 \mathbin{\times_X}\sqrt[N]{(X,C_s)}\longrightarrow X.
\]
Locally at a crossing this has the quotient chart
\[
 \left[\Spec R[u,v]/(u^N-x,v^N-y)\,/\,\mu_N^2\right],
\]
with the analogous one-variable chart away from crossings.
These charts are regular, and the stack is a regular tame
Deligne--Mumford stack, proper over $A$, with coarse space $X$.
Let $\mathfrak C_i$ be its reduced divisor above $C_i$. Then
$N\mathfrak C_i=\pi^*C_i$.

The sheaf $M$ extends to a finite locally constant sheaf on
$\mathfrak X_N$. Indeed the displayed root charts kill its
codimension-one tame inertia by the DVR form of Abhyankar's lemma.
Normalization on the regular charts is then \'etale by purity;
the canonical normalizations and their descent data agree on
chart overlaps. This proves the asserted extension, using only
Abhyankar's lemma and regular-scheme purity
\cite[tags 0BRM and 0BMB]{StacksAbhyankar,StacksBranchPurity}.
The same argument applies if $M$ has wild monodromy away from
these inertia groups.

The component $\mathfrak C_i$ is a $\mu_N$-gerbe over the root
curve obtained from $C_i$ by inserting $N$th roots at its $\nu_i$
crossing points. Degrees on $\mathfrak C_i$ will be normalized
so that the pullback of a line bundle on $C_i$ has its ordinary
degree. Concretely every line bundle has its $N$th power descending
to $C_i$, and we divide that descended degree by $N$. This convention
is additive and gives
\[
 \deg\bigl(\mathcal O(\mathfrak C_j)|_{\mathfrak C_i}\bigr)
        =B_{ij}/N,
 \qquad \deg(J|_{\mathfrak C_i})=d_i.
 \tag{F8.2}
\]
For vector bundles use the degree of the determinant. Nonzero
maps of line bundles do not increase degree, since their zeros
are effective divisors of nonnegative rational degree.
The degree multiplies by $p$ under absolute Frobenius, since
Frobenius pulls a line bundle back to its $p$th power.
The local charts and degree convention agree with the usual
root-stack/parabolic description
\cite[\S\S2.3--2.4]{Cadman2007}
\cite[Theorems 4--5, pp.~12,17]{Borne2007}.

Set $V=M\otimes_{\mathbf F_p}\mathcal O_{\mathfrak X_N}$, of rank $r$.
It has a unit Frobenius structure $F^*V\simeq V$ and is trivialized
by a finite \'etale cover. Consequently $V|_{\mathfrak C_i}$ is
semistable of degree zero: pull a subbundle back to that cover,
take its determinant inside an exterior power of a trivial bundle,
and use the preceding degree inequality. Under a finite representable
cover the normalized degrees scale by a positive rational factor:
ordinary stack degree scales by the degree of the cover, while
normalizing multiplies by the generic stabilizer order. Thus degree
signs are preserved. Its determinant has degree zero because it
becomes trivial on the cover. No restriction
on the degree of the trivializing cover is used here. For root curves with trivial generic stabilizer this is also
the finite-bundle semistability assertion of
\cite[Proposition 6, p.~22]{Borne2007}.

\paragraph{2. Coarse cohomology and a useful curve estimate.}
For a vector bundle $W$ on $\mathfrak C_i$, decompose it by the
characters of the generic $\mu_N$. All nontrivial generic-character
summands have zero coherent cohomology: taking invariants along
that gerbe is exact and kills them. The character-zero summand
$W_0$ descends to the root curve of $C_i$ at the crossings.
Its pushforward $E_0$ to $C_i$ is a vector bundle of the same rank.
On a local root chart a graded basis gives
\[
 0\leq \deg W_0-\deg E_0\leq\nu_i\operatorname{rank}W_0.
 \tag{F8.3}
\]
For clarity, adjunction from the pullback of $E_0$ is a generically
invertible map to $W_0$. In a character basis at a stacky point,
its matrix has diagonal entries $z^{a_j}$ with $0\leq a_j<N$.
The determinant degree lost in pushforward is $\sum_j a_j/N$.
Summing over the crossings proves (F8.3). Graded bases are obtained
by lifting a character basis modulo the maximal ideal and using
Nakayama's lemma; averaging makes the lifts equivariant.

Thus if $W_0$ is semistable of negative slope $t$, then it has no
sections and ordinary Riemann--Roch for $E_0$ gives
\[
 h^1(\mathfrak C_i,W_0)
 \leq \operatorname{rank}(W_0)(g_i-1-t+\nu_i).
 \tag{F8.4}
\]
If its slope is zero, evaluation at a nonstacky point is injective
on sections, since a section vanishing there gives a positive-degree
line subbundle after saturation. Hence
\[
 h^1(\mathfrak C_i,W_0)
 \leq \operatorname{rank}(W_0)(g_i+\nu_i).
 \tag{F8.5}
\]
These statements also apply to an isotypic summand tensored by a
line bundle that makes its generic character zero. Each isotypic
summand of a degree-zero semistable bundle is again semistable of
degree zero: it is a direct summand, every summand has degree at
most zero, and their degrees sum to zero.

Coarse pushforward on these tame stacks is exact on coherent
sheaves and on sheaves of $\mathbf F_p$-vector spaces. Locally it
is invariants for a finite group of order prime to $p$, followed
by finite affine pushforward; the invariants functor is exact by
averaging. For the \'etale assertion, take a strict henselian
neighborhood on the coarse space: its finite root cover is a finite
union of strictly henselian local schemes, whose \'etale global
section functors are exact. The quotient atlas then computes the
pushforward by finite-group cohomology, which vanishes in positive
degrees. The same description commutes with restriction to a closed
fibre. More explicitly, write $j:U\hookrightarrow\mathfrak X_N$ and
$j_X:U\hookrightarrow X$. Exact coarse pushforward gives
$R\pi_*Rj_*M=Rj_{X*}M$, and its closed-fibre compatibility gives
$i_X^*R\pi_*Rj_*M=R\pi_{D*}i^*Rj_*M$. These identities apply
to the bounded-below torsion complex $Rj_*M$ degree by degree.
Proper base change for the proper scheme $X$ over the strictly
henselian base $\Spec A$ therefore identifies
$H^1(U,M)$ with $\mathbb H^1(\mathfrak D_N,i^*Rj_*M)$.
Its spectral sequence yields the required nearby-cycle five-term
sequence. Coherent cohomology is computed by the same exact
coarse pushforward and the preceding coarse-space estimates.

\paragraph{3. A Frobenius-cokernel slope bound independent of $N$.}
Let
\[
 B_N=(F_*\mathcal O_{\mathfrak X_N})/\mathcal O_{\mathfrak X_N}.
\]
It is locally free: absolute Frobenius induces the automorphism
$\zeta\mapsto\zeta^p$ on each $\mu_N$ stabilizer, so it is
representable and is finite flat on the regular root charts. Its
unit section extends locally to a basis.
The pullback $F^*B_N$ is canonically isomorphic, as a module, to
the augmentation ideal in $F^*F_*\mathcal O$: the unit provides a
section of the augmentation. The powers of that ideal give a
filtration with locally free graded pieces
\[
 T^a(\Omega^1_{\mathfrak X_N/k}),\qquad 1\leq a\leq2(p-1),
 \tag{F8.6}
\]
where $T^a$ is the degree-$a$ truncated symmetric power.
This is Sun's canonical filtration
\cite[Theorem 3.7, pp.~12--13; local proof pp.~6--8]{Sun2008}.
In the present complete-local setting it can also be checked
directly after faithfully flat completion of a regular root chart:
with parameters $u,v$ and coefficient field $k$, the algebra is
\[
 k[[u,v]]\otimes_{k[[u^p,v^p]]}k[[u,v]]
 \simeq k[[u,v]][\alpha,\beta]/(\alpha^p,\beta^p),
\]
with augmentation ideal $(\alpha,\beta)$. Its graded generators
are $du,dv$. This calculation is invariant under changes of
parameters and therefore descends. All modules and maps involved are finite. Every closed point on
the proper exceptional divisor has residue field $k$, so its
completed regular chart is $k[[u,v]]$. Faithfully flat completion
checks these identities at every such point; any nonzero coherent
error supported on that divisor would have a closed point in its
support. Thus the calculation suffices on a neighborhood of the
exceptional divisor, which is all that the proof uses.

The conormal sequence on $\mathfrak C_i$ is an extension of two
line bundles, of degrees
\[
 -B_{ii}/N,\qquad
 2g_i-2+\nu_i(1-1/N).
 \tag{F8.7}
\]
The second formula follows from $dx=Nz^{N-1}dz$ at each root point;
the generic gerbe is \'etale and contributes no cotangent direction.
Both numbers are at most $K_i$. Each $T^a(\Omega^1)|_{\mathfrak C_i}$
has a filtration by the characters of this two-line flag: its
weights are $u,v$ with $0\leq u,v\leq p-1$ and $u+v=a$.
Thus its line-bundle layers have degrees at most $2(p-1)K_i$.
Tensoring with the semistable degree-zero bundle $V|_{\mathfrak C_i}$
shows that every subbundle of
$F^*(V\otimes B_N)|_{\mathfrak C_i}$ has slope at most
$2(p-1)K_i$. Here we used $F^*V\simeq V$ and the filtration (F8.6).
The assertion about subbundles follows by intersecting with the
filtration and adding degrees and ranks; each graded subbundle has
the indicated slope bound.

Pulling any subbundle of $(V\otimes B_N)|_{\mathfrak C_i}$ back by
Frobenius and dividing degrees by $p$ therefore proves the uniform
upper slope bound
\[
 \mu(\text{any subbundle of }(V\otimes B_N)|_{\mathfrak C_i})
       \leq b_i:=2(p-1)K_i/p.
 \tag{F8.8}
\]
This is independent both of $N$ and of $r$.

\paragraph{4. Nearby-cycle filtration and a uniform cutoff.}
Write $\mathfrak Z_N=\pi^*Z=\sum_iNc_i\mathfrak C_i$ and continue
to denote its ideal by $J$. The sheaf $M$ is locally constant on
$\mathfrak X_N$, so the Artin--Schreier calculation of F5/F7 applies
on its \'etale charts and descends. In particular, with
$\Psi=i^*j_*$ for the reduced exceptional stack,
\[
 0\to J^{-m}V/V\xrightarrow{F_V-1}J^{-pm}V/V
       \to U_mR^1\Psi M\to0,
 \tag{F8.9}
\]
and, for $m\geq2$,
\[
 \operatorname{gr}_mR^1\Psi M
  \simeq B_N|_{\mathfrak Z_N}\otimes V|_{\mathfrak Z_N}
                          \otimes J^{-m}|_{\mathfrak Z_N}.
 \tag{F8.10}
\]
The projection formula proves the latter identity, as in Suzuki's
Proposition 3.1; $F^{-1}\mathfrak Z_N=p\mathfrak Z_N$.

Filter $\mathcal O_{\mathfrak Z_N}$ by removing one irreducible
component at a time. A line-bundle layer $P$ on $\mathfrak C_i$
has degree bounded above by $L_i$, by (F8.2): all cycle coefficients
are between $0$ and $Nc_j$. The corresponding layer of (F8.10)
therefore has every subbundle of slope at most
\[
 b_i+L_i-md_i.
\]
Choose an integer $m_0\geq2$ strictly greater than
$\max_i(b_i+L_i)/d_i$. The last expression is negative for every
$m\geq m_0$, so there are no sections on any layer and hence
\[
 H^0(\operatorname{gr}_mR^1\Psi M)=0\quad(m\geq m_0).
 \tag{F8.11}
\]
The number of layers grows with $N$, but their common vanishing
bound does not. Quasi-compactness and the increasing union of the
pole lattices give
$H^0(R^1\Psi M)=H^0(U_{m_0}R^1\Psi M)$, increasing $m_0$ by one
if desired.

\paragraph{5. Low poles: cancellation of the growing number of layers.}
The negative-degree filtration of Suzuki's Proposition 2.1 works
on the root surface as well. Its proof only uses the negative
definiteness of the intersection matrix and the nonnegativity of
off-diagonal intersections; here the normalized matrix is $B/N$.
Thus $J^{-pm}/\mathcal O$ and $J^{-pm}/J^{-m}$ admit filtrations
by line bundles of strictly negative degree on the $\mathfrak C_i$.
Tensoring with $V$ gives no sections, because $V|_{\mathfrak C_i}$
is semistable of degree zero. Equation (F8.9) consequently injects
$H^0(U_mR^1\Psi M)$ into the kernel of a Frobenius-linear map minus
an injective linear map on
$H^1(J^{-m}V/V)$. The fixed-space argument from F5 bounds its
$\mathbf F_p$-dimension by $h^1(J^{-m}V/V)$.

We bound the latter dimension uniformly for the fixed $m=m_0$.
Choose the negative filtration of $J^{-m}/\mathcal O$ by adding
components to the cycle $m\mathfrak Z_N$. For a layer $P$ on
$\mathfrak C_i$ we have
\[
 -mL_i\leq\deg P<0.
 \tag{F8.12}
\]
Decompose $V|_{\mathfrak C_i}=\bigoplus_{\chi}V_{i,\chi}$ by its
generic $\mu_N$-characters, and let $r_{i,\chi}$ be their ranks.
Only the summand whose character cancels that of $P$ contributes
to cohomology. As the coefficient of $\mathfrak C_i$ increases
from $0$ to $mNc_i$, the generic character of the successive
layers runs through every residue modulo $N$ exactly $mc_i$
times. Contributions from other components have trivial generic
character at $\mathfrak C_i$. Therefore the sum of the ranks of
all contributing summands for this component is exactly
\[
 mc_i\sum_\chi r_{i,\chi}=mc_i r.
 \tag{F8.13}
\]
This is the needed cancellation of the extra $N$ layers.

Each contributing bundle $V_{i,\chi}\otimes P$ is semistable of
negative slope $\deg P$ and descends after killing the generic
gerbe. By (F8.4) and (F8.12), its first cohomology has dimension
at most its rank times $g_i+\nu_i+mL_i+1$. Summing the exact
filtration and using (F8.13) gives
\[
 h^1(J^{-m}V/V)
 \leq r\,m\sum_i c_i(g_i+\nu_i+mL_i+1).
 \tag{F8.14}
\]
The right side is independent of $N$. This bounds
$\dim_{\mathbf F_p}H^0(R^1\Psi M)$ by a constant times $r$.

\paragraph{6. The proper exceptional term.}
Let $\mathfrak D_N=\bigcup_i\mathfrak C_i$ with its reduced
structure. Its normalization sequence has one term at each crossing,
with stabilizer $\mu_N^2$. Invariants there have dimension at most
$r$ and no higher coherent cohomology. On a component, only the
generic-character-zero summand of $V$ contributes; it is semistable
of degree zero. Estimate (F8.5) thus yields
\[
 h^1(\mathfrak D_N,V)
 \leq r\left(\sum_i(g_i+\nu_i)+\#\{\text{crossings}\}\right).
 \tag{F8.15}
\]
The Artin--Schreier sequence identifies
$H^1(\mathfrak D_N,M)$ with the Frobenius-fixed subgroup of this
finite-dimensional coherent cohomology. Surjectivity of $F_V-1$
on global sections follows from the same finite-dimensional
semilinear argument as in F5. Its $\mathbf F_p$-dimension is
therefore at most (F8.15).

The nearby-cycle five-term sequence, justified by coarse descent
in step 2, combines (F8.14) and (F8.15) to prove (F8.0). One possible
constant is
\[
 C_X=m_0\sum_i c_i(g_i+\nu_i+m_0L_i+1)
            +\sum_i(g_i+\nu_i)+\#\{\text{crossings}\}.
\]
It depends on $X,J,p$ and not on the coefficient rank, the monodromy
group, or any tame inertia order.
\end{proof}

\begin{corollary}[F8b: all pro-$p$-by-pro-prime-to-$p$ quotients; Proved]
Every quotient $G$ of $\pi_1(U)$ admitting an
exact sequence $1\to P\to G\to\Gamma\to1$ with $P$ pro-$p$ and
$\Gamma$ pro-prime-to-$p$ is uniformly topologically finitely
generated. No abelian hypothesis on $\Gamma$ is necessary.
\end{corollary}
\begin{proof}
Every finite simple $\mathbf F_p[G]$-module kills $P$ and hence has
prime-to-$p$ monodromy. Such a system is tamely ramified on the
resolution. F8 and inflation give the uniform degree-one bound.
The quotient $\Gamma$ has at most a fixed number of generators,
because it is a quotient of $\pi_1(U)^{(p')}$, which is finitely
generated by Grothendieck--Murre. The relative criterion F7a gives
the conclusion. Alternating wild and tame normal factors remain
outside this assertion, as do general nonsolvable wild kernels.
\end{proof}

\subsection{Pro-$p$ extensions of the full tame quotient}

The tame finite-generation statement below follows by the standard
curve and gluing methods; the uniform extension consequence F9b
uses the coefficient estimate proved in F8. Both arguments and the
SGA~1 source passed an independent root audit. No historical
priority assertion is made.

\begin{lemma}[F9: a uniform generator bound for the tame quotient; Known method, proved here]
Keep the notation of F8, and let $\pi_1^t(U;X,D)$ denote the quotient
classifying finite covers tamely ramified along the fixed SNC divisor
$D$ on $X$. Let $\Delta$ be the connected dual graph of $D$. Then
\[
 d\bigl(\pi_1^t(U;X,D)\bigr)
 \leq d_X:=\sum_i(2g_i+\nu_i+1)+b_1(\Delta).
 \tag{F9.1}
\]
In particular this is a bound for the full tame group, not only its
maximal prime-to-$p$ quotient.
\end{lemma}
\begin{proof}
Let $N$ be prime to $p$, and use the root stack $\mathfrak X_N$
and reduced exceptional stack $\mathfrak D_N$ of F8. We prove
that every finite quotient of $\pi_1(\mathfrak X_N)$ satisfies
(F9.1), independently of $N$.

\emph{The closed fibre surjects on fundamental groups.}
Let $\mathfrak Y\to\mathfrak X_N$ be a connected finite \'etale
cover. Its coarse space $Y$ is a finite scheme over $X$.
Indeed, locally $\mathfrak X_N=[V/H]$ with $V$ finite over an
affine open $\Spec R\subset X$ and $H=\mu_N$ or $\mu_N^2$.
The pullback of the cover is $W\to V$ finite \'etale, with
$H$-equivariant descent, and its coarse space is
$\Spec\Gamma(W,\mathcal O_W)^H$. This invariant ring is finite
over $R$: it is an $R$-submodule of a finite $R$-module and $R$
is noetherian. These finite coarse schemes glue. Thus $Y$ is
proper over $A$, and it is connected because a coarse-space map
identifies the underlying topological spaces.

For a proper scheme over a henselian local ring, connectedness
is equivalent to connectedness of the closed fibre. For example,
apply the full faithfulness in
\cite[Lemma 58.9.1, tag 0A48]{StacksProperEtaleSpecialization}
to the scheme and its trivial two-sheeted cover: sections of that
cover correspond exactly to open-and-closed decompositions.
Consequently $Y_k$ is connected. Locally in the above quotient
charts, taking invariants commutes with the residue-field base
change by averaging, since $|H|$ is invertible. Therefore $Y_k$
is the coarse space of $\mathfrak Y_k$, which is connected.
Replacing the closed fibre by its reduced support does not change
its underlying topological space. It follows that the restriction
of every connected finite \'etale cover to $\mathfrak D_N$ stays
connected. This is precisely the surjectivity of
\[
 \pi_1(\mathfrak D_N)\longrightarrow\pi_1(\mathfrak X_N).
 \tag{F9.2}
\]
This argument requires only this surjectivity; it does not assume
a full stack version of the proper specialization equivalence.

\emph{A bound for each normalization component.}
The component $\mathfrak C_i$ is the banded $\mu_N$-gerbe over
the root curve $\mathcal C_i$ obtained from $C_i$ by taking roots
at its $\nu_i$ crossing points. In every finite quotient $G$ of
$\pi_1(\mathfrak C_i)$, the image $I$ of the generic relative
$\mu_N$ is cyclic and central. To see centrality, the band supplies
natural automorphisms of every fibre functor; naturality makes
them commute with its automorphisms. Quotienting the corresponding
torsor by $I$ kills the relative gerbe inertia, so the torsor
descends to a finite \'etale $G/I$-torsor on $\mathcal C_i$.
This descent can be checked on a neutralizing \'etale cover of the
gerbe, where it is exactly descent of an equivariant finite \'etale
algebra with trivial $\mu_N$-action.

Every finite \'etale cover of $\mathcal C_i$ restricts to a cover
of $C_i$ minus its crossing points, tamely ramified at those
points. The local chart is $[\Spec k[[z]]/\mu_N]$ with coarse
parameter $z^N$, so the inertia order divides $N$. A connected
cover remains connected on this dense open subset because the
root curve is normal. Hence
$\pi_1^t(C_i\setminus\{\text{crossings}\})$ surjects onto
$\pi_1(\mathcal C_i)$.
SGA~1, Expos\'e~XIII, Corollaire~2.12 gives $2g_i+\nu_i$
generators for this full tame curve group
\cite[original p.~392; retyped pp.~290--292]{SGA1TameCurves}.
It follows that every such $G$ needs at most
\[
                              2g_i+\nu_i+1               \tag{F9.3}
\]
generators: lift generators of $G/I$ and adjoin one generator of
$I$. The compactness criterion for generators gives the same bound
for $\pi_1(\mathfrak C_i)$ itself.

\emph{Gluing the components.}
Finite \'etale covers of $\mathfrak D_N$ are obtained by gluing
finite \'etale covers of its normalization components, with
identifications over the two preimages of each node gerbe.
The local assertion follows on the node chart
\[
 [\Spec k[[u,v]]/(uv)\,/\,\mu_N^2].
\]
The strictly henselian nodal local ring and its two strictly
henselian branches have only trivial finite \'etale covers. Thus
on the quotient chart all three cover categories are finite
$\mu_N^2$-sets, and the gluing statement is the identification
of these sets at the two branches. The assertion descends from
these charts and the smooth locus. Equivalently, it is the usual
normalization gluing of finite \'etale algebras, with equivariant
descent along a finite tame atlas.

For completeness, the resulting generator bound does not require
injectivity of any edge-group map. Choose a spanning tree of
$\Delta$, base points on the components, and paths to the incident
node points. The gluing identifications along the tree allow one
to use a single fibre. The monodromy is then generated by the
images of the component fundamental groups and one gluing
element for every edge outside the tree. If the subgroup they
generated in a finite connected torsor were proper, its orbits
would give a nontrivial decomposition of all component covers
compatible with every gluing identification, hence a decomposition
of that torsor. This contradicts connectedness. Thus the number
of generators is bounded by the sum of (F9.3) and $b_1(\Delta)$.
Together with (F9.2), this proves the asserted uniform bound for
every $\pi_1(\mathfrak X_N)$.

Every finite tame quotient of $\pi_1(U)$ extends to one of these
root stacks by the Abhyankar--purity argument of F8. Consequently
every finite quotient of $\pi_1^t(U;X,D)$ obeys the same bound.
The compactness argument for generating tuples now proves (F9.1).
\end{proof}

\begin{corollary}[F9b: a normal pro-$p$ kernel over a tame quotient; Proved]
Let $G$ be a quotient of $\pi_1(U)$, let $P\triangleleft G$ be
pro-$p$, and assume that the composite
\[
 \pi_1(U)\longrightarrow\Gamma:=G/P
\]
factors through $\pi_1^t(U;X,D)$. Then
\[
 d(G)\leq\max\{d_X,\lceil C_X\rceil+1\},
 \tag{F9.4}
\]
where $C_X$ is the uniform constant in F8. No solvability
hypothesis is imposed on $\Gamma$ or on its finite quotients.
\end{corollary}
\begin{proof}
Every finite simple $\mathbf F_p[G]$-module is killed by $P$,
because a normal finite $p$-group has nonzero invariants on every
nonzero representation in characteristic $p$, and the invariant
space is $G$-stable. The associated sheaf on $U$ therefore has
tame boundary monodromy. Inflation and F8 give the uniform
rank-linear bound on $H^1(G,M)$ for all these simple modules.
F9 bounds the number of generators of $\Gamma$. The relative
criterion F7a proves (F9.4).
\end{proof}

\paragraph{Remaining wild obstruction.}
The kernel of $\pi_1(U)\to\pi_1^t(U;X,D)$ is normally generated
by wild inertia groups. This does not make it a pro-$p$ group.
For example, a finite group generated by $p$-subgroups can have
minimal normal subgroups of another prime order: $S_3$ is generated
by its order-two subgroups and has normal subgroup $C_3$.
Thus F9b does not close the alternating tame/wild case or the full
local fundamental-group conjecture. No assertion of historical
priority is attached to F9 or its proof.

\subsection{Abelian wild inertia and metabelian quotients}

\begin{proposition}[F10: the abelian inertia kernel is finite; Proved]
Let $A,U,X,D=\bigcup_iC_i$ be as in F8. The kernel of
\[
 \pi_1(U)^{\mathrm{ab}}\longrightarrow\pi_1(X)^{\mathrm{ab}}
 \tag{F10.1}
\]
is finite. More explicitly, write $B=(C_i\cdot C_j)$ and set
\[
 k_i=\max\{0,2g_i-2+\nu_i\},\qquad
 t=(-B)^{-1}k.
 \tag{F10.2}
\]
Every $p$-primary rank-one character $\chi$ of $\pi_1(U)$ has
Swan conductor $s_i(\chi)\leq t_i$ at $C_i$.
Choose $e\geq0$ with $p^e>\max_i t_i$. Then the kernel in
(F10.1) is annihilated by $p^e|\det B|$.
\end{proposition}

\paragraph{Status and source.}
The complete argument passed an independent root proof and primary-source
audit; BCE separately passed F10 and F10b. The input used here is the
global refined Swan conductor on regular $F$-finite
schemes in Krishna--Majumder, arXiv:2601.12177v2, 8 November 2026,
Corollary 9.7 p.~56 (Theorem 9.4 pp.~55--56), with the hypotheses
on p.~2. The source is a preprint. Its formulation includes
complete-local resolutions and does not require algebraization
or finite type over a field. No priority assertion is made.

\begin{proof}
A coefficient field identifies $A$ as a quotient of a power-series
ring over perfect $k$, so $A$ and the resolution charts are
$F$-finite. The resolution is an equidimensional noetherian regular
$\mathbf F_p$-scheme with SNC boundary. These are precisely the
hypotheses of the cited global refined-conductor theorem.
All differential forms over $\mathbf F_p$ below may also be viewed
over $k$, since $k$ is perfect and $\Omega^1_{k/\mathbf F_p}=0$.

Fix a character $\chi\in H^1(U,\mathbf Q_p/\mathbf Z_p)$.
It has finite order, say dividing $p^m$. Let
\[
                         S=\sum_i s_i C_i
\]
be its logarithmic Swan-conductor divisor. Each $s_i$ is a
nonnegative integer, defined at the henselian discrete valuation
field at the generic point of $C_i$. If $s_i>0$, the character is
in $\operatorname{fil}^{\log}_S H^1(U,\mathbf Z/p^m)$ and is not
in $\operatorname{fil}^{\log}_{S-C_i}H^1(U,\mathbf Z/p^m)$.
Since $S-C_i\geq\lfloor S/p\rfloor$, Corollary 9.7 of
\cite{KrishnaMajumder2026} gives a nonzero section
\[
 0\ne\operatorname{Rsw}_i(\chi)\in
 H^0\!\left(C_i,
   \Omega_X^1(\log D)|_{C_i}\otimes\mathcal O_X(S)|_{C_i}\right).
 \tag{F10.3}
\]
The map is a levelwise monomorphism, not only an injection after
passing to an unspecified colimit; see Theorem 9.4(1).
Its generic stalk is the local refined Swan conductor by
Corollary 9.6. Thus the nonvanishing in (F10.3) has exactly the
required meaning.

The residue sequence on $C_i$ is
\[
 0\longrightarrow\Omega^1_{C_i/k}(\log\{\text{crossings}\})
 \longrightarrow\Omega_X^1(\log D)|_{C_i}
 \longrightarrow\mathcal O_{C_i}\longrightarrow0.
 \tag{F10.4}
\]
It follows either from the local coordinates $d\log x,dy$ and
$d\log x,d\log y$, or from Lemma 9.8 of the same source.
The two line bundles have degrees $2g_i-2+\nu_i$ and $0$.
Consequently every line subbundle of the middle term has degree
at most $k_i$. A nonzero section in (F10.3), after saturation,
therefore implies
\[
                              S\cdot C_i\geq-k_i.
 \tag{F10.5}
\]
When $s_i=0$, this inequality holds as well, because
$S\cdot C_i=\sum_{j\ne i}s_j(C_j\cdot C_i)\geq0$.
Thus, as a vector inequality, $Bs\geq-k$.

We include the elementary matrix step. Put $C=-B$. This matrix
is symmetric positive definite and has nonpositive off-diagonal
entries. Its inverse is entrywise nonnegative. Indeed, if
$Cx=v\geq0$, write $x=x_+-x_-$ with $x_+,x_-$ nonnegative and
of disjoint support. If $x_-\ne0$, then
\[
 0\leq x_-^Tv=x_-^TCx_+-x_-^TCx_-<0,
\]
because the mixed term is nonpositive and the quadratic term is
positive. This contradiction proves $x\geq0$.
Applying $C^{-1}$ to $Cs\leq k$ now gives $s\leq C^{-1}k=t$.
This proves the uniform conductor bound.

\emph{A fixed $p$-power kills all abelian wild inertia.}
We recall explicitly why multiplication by $p$ reduces the bound
by a factor $p$. For a henselian characteristic-$p$ DVR $R$ with
fraction field $L$, the Brylinski filtration on Witt vectors
$a=(a_{m-1},\ldots,a_0)\in W_m(L)$ is
\[
 \operatorname{fil}_nW_m(L)
   =\{a:p^jv(a_j)\geq-n\text{ for all }j\}.
\]
Its image under Artin--Schreier--Witt is the logarithmic conductor
filtration. Witt restriction drops the last displayed coordinate
and gives
\[
 R:\operatorname{fil}_nW_m(L)
       \longrightarrow\operatorname{fil}_{\lfloor n/p\rfloor}
                         W_{m-1}(L).
 \tag{F10.6}
\]
This is Lemma 3.2, p.~12, of \cite{KrishnaMajumder2026}; the
classical characteristic-$p$ description is also given in
\cite[\S3.2.2--3.2.4, pp.~14--15]{KatoLealSaito2019}.
Under the standard embeddings $\mathbf Z/p^m\hookrightarrow
\mathbf Q_p/\mathbf Z_p$, restriction represents multiplication
by $p$. Hence
\[
 s_i(p^e\chi)\leq\lfloor s_i(\chi)/p^e\rfloor=0.
\]
Conductor zero here means unramified: its Witt representative is
integral, and the Artin--Schreier--Witt equation then defines an
\'etale cover of the DVR. Thus $p^e\chi$ is unramified along every
$C_i$. Purity on the regular $X$ extends it to $X$.

\emph{The prime-to-$p$ part.}
The Kummer/intersection calculation in F5 shows that multiplying
any prime-to-$p$ character by $b=|\det B|$ kills all its inertia
along the $C_i$. Purity again extends the resulting character
to $X$. One can see the calculation directly: extend its torsion
line bundle to $X$, write the divisor of the trivialization of
its $n$th power as $\sum a_i C_i$, and intersect with $C_j$.
Then $Ba\equiv0\pmod n$, so $ba_i\equiv0\pmod n$ by the
adjugate matrix. These residues describe the tame inertia.

Set $Q=\pi_1(U)^{\mathrm{ab}}$ and let $I$ be the kernel in
(F10.1). The map is surjective, since a connected finite \'etale
cover of normal $X$ stays connected on its dense open subset $U$.
The previous two paragraphs show that all characters of $I$
obtained by restriction from $Q$ vanish on $bp^eI$. Finite
characters of the profinite abelian group $Q$ separate its
points, so $bp^eI=0$. By Grothendieck--Murre and Suzuki, $Q$ is
topologically finitely generated. Its subgroup of elements
annihilated by the fixed integer $bp^e$ is finite. Therefore $I$
is finite, as asserted.
\end{proof}

\begin{corollary}[F10b: one fixed cover removes every abelian inertia action; Proved]
There is a finite connected abelian cover $U'\to U$ such that,
for every finite locally constant sheaf on $U$ with abelian
monodromy, its pullback to $U'$ extends to a finite locally constant
sheaf on a fixed SNC resolution $X'$ of the normalization of $A$
in $U'$. The same cover and resolution work for all coefficient
primes and all ranks.
\end{corollary}
\begin{proof}
Choose an open subgroup $Q'\subset Q$ with $Q'\cap I=0$; this
is possible because $I$ is finite and $Q$ is profinite. Let $U'$
be the cover corresponding to the inverse image of $Q'$.
Its normalization $A'$ is finite over $A$, excellent, normal,
complete and local with residue field $k$. Fix an SNC resolution
$X'$ of $A'$. The pulled-back abelian monodromy factors through
$Q'$.

At a height-one exceptional point of $X'$, let $R$ be its DVR.
The morphism $\Spec\operatorname{Frac}(R)\to X$ extends to
$\Spec R\to X$ by properness of $X$ over $A$. Every finite
character of $Q/I$, being a cover of $X$, is unramified after this
pullback. Hence the image in $Q$ of the inertia at $R$ lies in
$I$. It also lies in $Q'$, so it is trivial. Thus the given pulled-back
monodromy is unramified at every such DVR. It is already \'etale
elsewhere. Normalization and regular-scheme purity now extend the
corresponding finite cover to $X'$. This proof uses one fixed
$U'$ and $X'$, independently of the sheaf.
\end{proof}

\begin{proposition}[F10c: all-prime bounds for abelian coefficient monodromy; Proved]
There is a constant $C_A$ such that for every prime $\ell$ and
every finite locally constant $\mathbf F_\ell$-sheaf $M$ on $U$
with abelian monodromy,
\[
                       \dim_{\mathbf F_\ell}H^1(U,M)
                       \leq C_A\operatorname{rank}M.
 \tag{F10.7}
\]
\end{proposition}
\begin{proof}
Use the fixed cover and resolution of F10b. If $\ell=p$, F7
bounds the cohomology after pullback by $C_{X'}\operatorname{rank}M$.
If $\ell\ne p$, every degree-one cocycle after pullback already
factors through the tame quotient. Indeed, it gives an affine
monodromy representation into $M\rtimes\operatorname{im}\rho$.
The original representation $\rho$ is unramified on $X'$, so its
restriction to boundary inertia is trivial. On inertia the cocycle
is a homomorphism to the additive $\ell$-group $M$. It kills wild
pro-$p$ inertia, and the affine cover is therefore tame.
Consequently
\[
 H^1(U',M)=H^1(\pi_1^t(U';X',D'),M),
\]
whose dimension is at most $d_{X'}\operatorname{rank}M$ by F9
and the values of a cocycle on generators.

Let $F=\operatorname{Gal}(U'/U)$. Inflation--restriction gives
\[
 0\to H^1(F,M^{\pi_1(U')})\to H^1(U,M)
          \to H^1(U',M)^F.
\]
The first term has dimension at most $d(F)\operatorname{rank}M$,
with no assumption that $|F|$ is prime to $\ell$.
Thus one can take
$C_A=d(F)+\max\{C_{X'},d_{X'}\}$.
\end{proof}

\begin{theorem}[F11: all metabelian local quotients are finitely generated; Proved]
Every metabelian quotient $G$ of $\pi_1(U)$ is topologically
finitely generated, with a bound depending only on $A$.
Both tame and wild normal factors are allowed.
\end{theorem}
\begin{proof}
Let $d_0=d(\pi_1(U)^{\mathrm{ab}})$ and let $C_A$ be as in
F10c. Put $D=\max\{d_0,\lceil C_A\rceil+1\}$.
We show that every finite metabelian quotient $F$ is generated
by $D$ elements. If $F'=1$, the bound $d_0$ applies. Otherwise
choose a minimal nonzero $F$-normal subgroup $M\subset F'$.
It is an elementary abelian $\ell$-group for some prime $\ell$,
and is simple as an $\mathbf F_\ell[F]$-module. Since $F'$ is
abelian, conjugation on $M$ factors through the abelian group
$F/F'$. Inflation from $F$ to $\pi_1(U)$ and F10c give
\[
                  \dim H^1(F,M)\leq C_A\dim M.
 \tag{F11.1}
\]
Inductively choose $D$ generators of $F/M$ and consider their
$|M|^D$ possible lifts. A nongenerating lift spans a proper
subgroup $J$ with $JM=F$. The subgroup $J\cap M$ is $F$-stable:
$J$ preserves it, and $M$ centralizes it. Minimality therefore
forces $J\cap M=0$, so $J$ is a complement. If there is no
complement, all lifts generate.

In the split case write $H=F/M$. Complements are indexed by
$Z^1(H,M)$. Since $M$ acts trivially on itself, splitting gives
\[
 H^1(F,M)\simeq H^1(H,M)\oplus\End_H(M).
\]
The endomorphism ring contains the scalar field $\mathbf F_\ell$.
Moreover $|B^1(H,M)|\leq|M|$. Hence (F11.1) implies
\[
 |Z^1(H,M)|\leq |M|^{C_A+1}/\ell<|M|^D.
\]
Each complement contains at most one lift of the chosen generating
tuple of $H$, so some lift generates $F$. This proves the bound
by induction. Finally the compactness of generating-tuple conditions
on $G^D$ gives $D$ topological generators of $G$.
\end{proof}

\paragraph{Scope remaining after these results.}
F11 closes derived length two, including the metabelian
abstract obstructions of F3 in every residue characteristic.
It does not treat arbitrary solvable quotients of larger derived
length: their minimal normal modules can have nonabelian
coefficient monodromy. Nor does it settle arbitrary nonsolvable
quotients. These are genuine remaining gates.

\subsection{The coefficient-$p$ half at derived length three}

Throughout this subsection, $A$ is the complete normal two-dimensional
local domain of equal characteristic $p>0$ with algebraically closed
residue field used in F8--F11, $U$ is its punctured spectrum, and
all resolutions are fixed SNC resolutions. No mixed-characteristic
assertion is made here. The proof passed an independent root audit.

\begin{lemma}[F12a: normal $p$-subgroups on simple characteristic-$p$ modules]
Let $\Gamma$ be a finite group and let $V$ be a nonzero simple
$\mathbf F_p[\Gamma]$-module. Every normal $p$-subgroup of
$\Gamma$ acts trivially on $V$. Consequently, if $\Gamma$ is
metabelian and its action on $V$ is faithful, then $\Gamma'$ has
order prime to $p$.
\end{lemma}
\begin{proof}
Every nonzero finite-dimensional representation of a finite
$p$-group $P$ over $\mathbf F_p$ has nonzero invariants. Here is an
inductive proof. A central subgroup $C_p\subset P$ exists; its
generator $u$ satisfies $(u-1)^p=0$, so the kernel of $u-1$ is
nonzero and $P$-stable. The quotient group $P/C_p$ acts on it,
and induction gives a nonzero $P$-fixed vector. If $P$ is normal
in $\Gamma$, the invariant space $V^P$ is $\Gamma$-stable, so
simplicity gives $V^P=V$.
For metabelian $\Gamma$, the $p$-primary part of its abelian
commutator group is characteristic in $\Gamma'$ and therefore
normal in $\Gamma$. It acts trivially; faithfulness makes it
trivial.
\end{proof}

\begin{proposition}[F12: characteristic-$p$ coefficients with metabelian monodromy; Proved]
There is a constant $C_{A,2,p}$ such that every simple finite
$\mathbf F_p$-local system $M$ on $U$ whose monodromy group is
metabelian satisfies
\[
                    \dim_{\mathbf F_p}H^1(U,M)
                    \leq C_{A,2,p}\operatorname{rank}M.
 \tag{F12.1}
\]
The system itself may be wildly ramified on the original resolution.
\end{proposition}
\begin{proof}
Use once and for all the finite connected abelian cover
$U_0\to U$ and the fixed resolution $X_0$ supplied by F10b.
Write $F=\operatorname{Gal}(U_0/U)$ and let
$\rho:\pi_1(U)\twoheadrightarrow\Gamma\subset\operatorname{GL}(M)$ be the
faithful monodromy image. It is metabelian by hypothesis.
Lemma F12a makes $\Gamma'$ a prime-to-$p$ group.

The finite quotient $\pi_1(U)\to\Gamma^{\mathrm{ab}}$ is an
abelian local system. Its pullback to $U_0$ extends to $X_0$,
by the uniform assertion of F10b. Thus at every exceptional
DVR of $X_0$, the image under $\rho$ of boundary inertia lies
in $\Gamma'$. Wild inertia is pro-$p$, so its image in the
prime-to-$p$ group $\Gamma'$ is trivial. Therefore the entire
pulled-back system $M|_{U_0}$ is tame along the fixed exceptional
divisor of $X_0$.

No simplicity assumption on this restricted system is needed:
F8 applies to arbitrary finite $\mathbf F_p$-local systems with
tame boundary monodromy. Hence
\[
 \dim H^1(U_0,M)\leq C_{X_0}\operatorname{rank}M.
\]
Inflation--restriction bounds the kernel of restriction by
$\dim H^1(F,M^{\pi_1(U_0)})\leq d(F)\operatorname{rank}M$.
Thus $C_{A,2,p}=C_{X_0}+d(F)$ works. In particular this argument
does not assume that $[U_0:U]$ is prime to $p$ or that restriction
preserves simplicity.
\end{proof}

\begin{corollary}[F12b: a pro-$p$ second derived subgroup; Proved]
Every quotient $G$ of $\pi_1(U)$ whose closed second derived
subgroup $G''$ is pro-$p$ is topologically finitely generated,
with a bound depending only on $A$.
\end{corollary}
\begin{proof}
The quotient $\Gamma=G/G''$ is metabelian, so F11 bounds its
number of generators uniformly. Every finite simple
$\mathbf F_p[G]$-module kills the normal pro-$p$ group $G''$ by
F12a applied to its finite image. Its faithful monodromy group
is therefore metabelian, so inflation and F12 give the uniform
$H^1$ bound. The relative criterion F7a proves the assertion.
\end{proof}

\paragraph{Precise remaining derived-length-three gate.}
For a finite group of derived length at most three, a minimal
normal subgroup inside its second derived group is an elementary
abelian $\ell$-group. Its conjugation representation is simple
and factors through a metabelian group, because the second derived
group is abelian. F12 supplies the needed uniform estimate when
$\ell=p$. The cases $\ell\ne p$, with possibly wild metabelian
coefficient monodromy, remain open here. F12b is a consequence
for a specified class, not a claim that every second derived group
is pro-$p$ or that bounded-exponent inertia inside a finitely
generated metabelian group is finite.

\subsection{A fixed logarithmic conductor bound for metabelian monodromy}

Throughout, $A$ is a complete normal two-dimensional local domain
of equal characteristic $p>0$ with algebraically closed residue field,
and $U$ is its punctured spectrum. Fix the cover $U_0\to U$ and
SNC resolution $X_0$ of its normalization supplied by F10b. Write
$D_0=\bigcup_i C_i$, $B=(C_i\cdot C_j)$,
\[
 k_i=\max\{0,2g(C_i)-2+\nu_i\},\qquad
 t=(-B)^{-1}k,\qquad T_0=\max_i t_i,
 \tag{F13.1}
\]
where $\nu_i$ is the number of nodes on $C_i$. As in F10, choose
a further blowup if necessary so that $D_0$ is nonempty.
The logarithmic breaks below are computed over the completed DVRs
at the generic points of $D_0$. The complete proof passed an
independent root audit; F13b also passed a separate BCE geometry audit.

\begin{lemma}[F13a: modular monomial representations with a specified normal subgroup]
Let $\Gamma$ be a finite group, let $N\triangleleft\Gamma$ be
abelian with $\Gamma/N$ abelian, and let $\kappa$ be an algebraically
closed field of characteristic $\ell$. Every nonzero simple
$\kappa[\Gamma]$-module is induced from a one-dimensional character
of a subgroup $K$ containing $N$. In particular $K$ is normal in
$\Gamma$ and $\Gamma/K$ is abelian. No restriction on the
$\ell$-part of $|\Gamma|$ is imposed.
\end{lemma}
\begin{proof}
The $\ell$-primary subgroup $N_\ell$ is normal in $\Gamma$ and
acts trivially by the normal-primary invariants argument of F12a,
with $p$ replaced by $\ell$. We may quotient it out. Restriction
to the resulting prime-to-$\ell$ abelian group $N$ is a direct
sum of character spaces. Simplicity makes the characters occurring
in $V$ a single $\Gamma$-orbit. Choose one character $\chi$, let
$I$ be its stabilizer, and put $W=V_\chi$. The character-space
sum gives $V=\operatorname{Ind}_I^\Gamma W$. Moreover $W$ is
simple as an $I$-module: a proper nonzero submodule would induce
a proper nonzero submodule of $V$.

Put $A_1=I/N$. The pairing
\[
 b:A_1\times A_1\longrightarrow\kappa^*,\qquad
 b(\bar x,\bar y)=\chi([x,y])
 \tag{F13.2}
\]
is well-defined, alternating and bimultiplicative. Indeed,
$[I,I]\subset N$, and the $I$-invariance of $\chi$ kills
$\chi([x,n])$ for $n\in N$; the commutator identities then
prove independence of lifts and bimultiplicativity. Its radical
contains the $\ell$-primary part of $A_1$, since $\kappa^*$ has
no nontrivial $\ell$-power torsion. Choose a maximal isotropic
subgroup $L\subset A_1$. It contains the radical and satisfies
$L=L^\perp$: an element of $L^\perp\setminus L$ could otherwise
be adjoined to produce a larger isotropic subgroup.

Let $K$ be the inverse image of $L$ in $I$. Commutators of
$K$ act on $W$ through the scalar $\chi([K,K])=1$, so the
finite family of matrices from $K$ commutes. Over $\kappa$ it
has a common eigenvector, defining a character
$\theta:K\to\kappa^*$ extending $\chi$. Frobenius reciprocity
gives a nonzero map
$\operatorname{Ind}_K^I\theta\to W$.
The stabilizer of $\theta$ in $I$ is $K$: for $x\in I$, equality
of $\theta$ and its conjugate means precisely
$b(\bar x,L)=1$, hence $\bar x\in L^\perp=L$.
Thus the restriction of $\operatorname{Ind}_K^I\theta$ to
$K$ is a direct sum of pairwise distinct character lines.
Every submodule is a sum of those lines, and $I$ acts transitively
on them. The induced module is therefore simple and the nonzero
map is an isomorphism. Induction in stages gives
$V=\operatorname{Ind}_K^\Gamma\theta$.
The preimage of $K$ before quotienting $N_\ell$ contains the
original $N$, as required. Finally every subgroup containing $N$
is normal in $\Gamma$, because $\Gamma/N$ is abelian.
\end{proof}

\begin{lemma}[F13b: the Swan bound is unchanged by a finite \'etale resolution cover]
Let $f:X_1\to X_0$ be any connected finite \'etale cover, and put
$U_1=f^{-1}(U_0)$ and $D_1=f^{-1}(D_0)$. If $C'_a$ lies over
$C_i$ with degree $d_a$, every finite $p$-primary character on
$U_1$ has Kato Swan conductor at $C'_a$ at most $t_i$.
In particular the bound is independent of $\deg f$.
\end{lemma}
\begin{proof}
The surface $X_1$ is regular and proper over $A_0$, the
normalization underlying $U_0$. Its Stein factor is the spectrum
of the normalization of $A_0$ in the connected cover $U_1$.
This normalization is finite, complete and local with residue
field $k$; $X_1$ is a resolution of it, and $D_1$ is its reduced
exceptional SNC divisor. These assertions can also be checked
from the finite \'etale morphism over $X_0$: away from $D_0$
it is exactly the specified finite \'etale cover, regularity
and the SNC condition are preserved, and the connected proper
closed fiber follows from proper henselian specialization.
Thus the proof of F10 applies to $X_1$.

Let $B'$ be the intersection matrix of the components $C'_a$.
The finite \'etale curve map $C'_a\to C_i$ gives
\[
 2g(C'_a)-2=d_a(2g(C_i)-2),\qquad \nu'_a=d_a\nu_i.
 \tag{F13.3}
\]
Consequently $k'_a=d_a k_i$. The projection formula gives, for
any rational vector $v=(v_i)$,
\[
 \left(B'f^*v\right)_a
 =\left(f^*\sum_i v_iC_i\right)\cdot C'_a
 =d_a(Bv)_i,
 \qquad (f^*v)_a=v_i.
 \tag{F13.4}
\]
Taking $v=t$ shows $(-B')f^*t=k'$, so
$(-B')^{-1}k'=f^*t$. Apply F10's refined-Swan inequality
$B's'\geq-k'$ and its nonnegative inverse-matrix argument:
$s'_a\leq(f^*t)_a=t_i$. This uses the same
Krishna--Majumder levelwise refined-Swan input as F10; it does
not assert a bound on the total order of the character.
\end{proof}

\begin{proposition}[F13: a uniform logarithmic conductor after one fixed cover; Proved]
For every prime $\ell\ne p$ and every finite locally constant
$\mathbf F_\ell$-sheaf $M$ on $U$ with metabelian monodromy,
the highest logarithmic break of $M|_{U_0}$ at $C_i$ is at most
$t_i$. The fixed cover and all the bounds are independent of
$\ell$, the rank and the finite metabelian monodromy group.
\end{proposition}
\begin{proof}
Let $\Gamma$ be the monodromy group of $M$, and let
$\Gamma_0$ be the image of $\pi_1(U_0)$. Put
$N=\Gamma_0\cap\Gamma'$. It is abelian and normal, and
$\Gamma_0/N$ is abelian. By F10b, the pullback of the original
abelian quotient $\Gamma^{\mathrm{ab}}$ extends over $X_0$.
Therefore the image of every boundary inertia group in
$\Gamma_0$ lies in this specified subgroup $N$.

Extend scalars to $\overline{\mathbf F}_\ell$ and take a
composition series of the restricted representation. For a
simple constituent $V$, F13a gives
$V=\operatorname{Ind}_K^{\Gamma_0}\theta$ with $K\supset N$.
The associated cover $U_1\to U_0$ is unramified at every
exceptional DVR, because all boundary inertia lies in $N$.
Purity on the regular surface $X_0$ extends its normalization
to a finite \'etale cover $f:X_1\to X_0$.
The character $\theta$ is a rank-one local system on $U_1$.
All its values lie in a finite extension of $\mathbf F_\ell$.
Split $\theta$ into its $p$-primary part and its part of order
prime to $p$. The latter is tame and has logarithmic break zero.
The former has Kato Swan conductor at $C'_a$ at most $t_i$
by F13b.

Kato--Saito, Theorem~1.3, identifies the rank-one Kato Swan
conductor with the Abbes--Saito logarithmic conductor over the
completed DVR, including imperfect residue fields
\cite[Theorem~1.3, p.~2]{KatoSaito2019}.
Identifying the finite cyclic character values with complex
roots of unity does not change its kernel or either conductor.
Thus the logarithmic break of $\theta$ at $C'_a$ is at most
$t_i$.

For precision, induction through this finite \'etale cover does
not multiply the break by its degree. At the completed DVR of
$C_i$, each local field extension belonging to a component of
$f$ is unramified. For any positive rational $r$, Abbes--Saito
base-change compatibility identifies its logarithmic group
$G^r_{\log}$ with the same subgroup of the original absolute
Galois group \cite[Proposition~3.15(3)]{AbbesSaito2002}.
For $r>t_i$, this subgroup acts trivially on each conjugate
character in the restricted induced representation, and trivially
on the cosets of every unramified local extension. It therefore
acts trivially on the induced representation. This proves the
claim for each simple constituent.

Finally the image of any positive logarithmic ramification group
is a finite $p$-group. Its representations in characteristic
$\ell\ne p$ are semisimple. Hence triviality on every constituent
implies triviality on the whole representation. Scalar extension
is faithful, so the same assertion holds for $M$. This last step
also shows that no semisimplicity of $M$ was assumed.
\end{proof}

\paragraph{The remaining complete-local Betti gate.}
F13 is a conductor estimate, not yet an $H^1$ estimate. The precise
additional assertion that would close the coefficient-$\ell\ne p$
case at derived length three is: for this fixed complete-local
pair $(X_0,D_0)$ and the fixed bounds $t_i$, there is a constant
$C$ such that
\[
 \dim_\Lambda H^1(U_0,L)\leq C\operatorname{rank}_\Lambda L
 \tag{F13.5}
\]
for all finite fields $\Lambda$ of characteristic different from
$p$ and all finite local systems $L$ whose logarithmic breaks
at $C_i$ are bounded by $t_i$. It would suffice to establish
(F13.5) just for the systems occurring in F13.
Inflation--restriction for the fixed $U_0/U$ and the complement
count of F11 would then give a uniform generator bound for all
quotients of derived length at most three. No such conclusion is
being made without (F13.5).

Hu--Teyssier prove the corresponding rank-linear estimate for a
proper smooth \emph{finite-type} $k$-scheme $\overline V$ and
$V=\overline V\setminus E$, uniformly in finite coefficient fields
\cite[Theorem~7.34, p.~41]{HuTeyssier2025Betti}. Their relative
Theorem~7.15 likewise requires finite-type schemes
(Definition~7.14 and Theorem~7.15, p.~33). The nearby-cycle
Theorem~4.7 requires a geometric henselian trait
(Definition~4.6 and Theorem~4.7, pp.~15--16), and its proof uses
that condition in the spreading-out step. Neither statement,
as printed, applies to the fixed complete-local pair above.
A comparison, approximation or local Betti theorem with constants
uniform for this pair remains necessary. This explicit scope
check prevents treating finite-type spreading out for one sheaf
as a uniform result for all ranks and all sheaves.


\subsection{Removing unramified order from the remaining rank-one problem}

The following reduction does not assert the complete-local Betti bound
(F13.5). It identifies a bounded-degree geometric problem sufficient for
that bound on the metabelian systems actually used in F13. The
complete reduction passed an independent root proof audit.

\begin{lemma}[F13c: bounded ramified order after an unramified twist]
Retain $X_0,D_0,U_0$ and $t_i$ from F13, and choose $e\geq0$
such that $p^e>\max_i t_i$. Let $f:X_1\to X_0$ be any connected
finite \'etale cover, with $U_1=f^{-1}(U_0)$. Every character
\[
 \chi\in H^1(U_1,\mathbf Q_p/\mathbf Z_p)
\]
can be written $\chi=\alpha+\psi|_{U_1}$, where
$p^e\alpha=0$ and $\psi\in H^1(X_1,\mathbf Q_p/\mathbf Z_p)$.
The same $e$ works for every such cover. The Swan conductor of
$\alpha$ at every exceptional component equals that of $\chi$.
\end{lemma}
\begin{proof}
First, $H^1(X_1,\mathbf Q_p/\mathbf Z_p)$ is divisible. Indeed,
proper henselian specialization identifies the finite \'etale cover
categories of $X_1$ and its reduced proper nodal closed fibre $D_1$
\cite[Lemma 58.9.1, tag 0A48]{StacksProperEtaleSpecialization}.
Thus their groups of finite continuous characters agree. To verify
divisibility without any assertion about the nonabelian fundamental
group, use the Artin--Schreier sequence on $D_1$. The group
$H^1(D_1,\mathcal O_{D_1})$ is a finite-dimensional $k$-vector space,
and $H^2(D_1,\mathcal O_{D_1})=0$. On this vector space, Frobenius
is additive and $p$-semilinear. The map $F-1$ is therefore a
homomorphism of vector algebraic groups with differential $-1$.
It is \'etale and its image is an open subgroup of a connected
vector group, hence is the whole group. Since $k$ is algebraically
closed, it is surjective on $k$-points. Artin--Schreier now gives
$H^2(D_1,\mathbf F_p)=0$.

The exact sequence
\[
 0\longrightarrow\mathbf F_p\longrightarrow
 \mathbf Q_p/\mathbf Z_p\xrightarrow{p}
 \mathbf Q_p/\mathbf Z_p\longrightarrow0
\]
shows that multiplication by $p$ is surjective on
$H^1(D_1,\mathbf Q_p/\mathbf Z_p)$. All of its elements are finite
characters, and the same is true on $X_1$. This proves the required
divisibility.

The F13b conductor bound applies to every rank-one character on
$U_1$, not only to the systems in the statement of F13.
By that bound and the Witt-restriction calculation (F10.6), the character
$p^e\chi$ has conductor zero along every exceptional component.
It extends to $X_1$ by purity. Choose an unramified character
$\psi$ whose $p^e$ multiple is this extension, using divisibility.
Then $\alpha=\chi-\psi|_{U_1}$ is killed by $p^e$. An unramified
twist does not change a positive logarithmic break and preserves
conductor zero; hence the Swan conductors are unchanged.
\end{proof}

\begin{proposition}[F13d: a bounded-degree sufficient criterion]
Keep the notation and $e$ of F13c. Suppose that there is a constant
$C$ with the following property. For every connected finite \'etale
cover $X_1\to X_0$ of degree $d$ and every connected cyclic cover
$V\to U_1$ of degree dividing $p^e$ corresponding to a character
whose Swan conductors satisfy the F13b bounds, the normalization
of the complete local surface underlying $X_1$ in $V$ admits an
SNC resolution $Y$ dominating $X_1$ such that
\[
 d\bigl(\pi_1^t(V;Y,D_Y)\bigr)\leq Cd.
 \tag{F13.6}
\]
Then there is a uniform rank-linear $H^1(U_0,-)$ bound for every
coefficient-$\ell\ne p$ system occurring in F13. Consequently
(F13.6), together with the already proved F12, would settle the
uniform finite-generation assertion for quotients of derived
length at most three.
\end{proposition}
\begin{proof}
First consider a rank-one character $\theta$ on $U_1$ with finite
values in a field of characteristic $\ell\ne p$. Split its finite
image into its $p$-primary and prime-to-$p$ parts. By F13c, after
a finite extension of the coefficient field if necessary, write
\[
                         \theta=\alpha\otimes\beta,
\]
where $\alpha$ has order dividing $p^e$ and $\beta$ is tame
along $D_1$. Here the $p$-primary factor removed from $\alpha$
extends over $X_1$, while every prime-to-$p$ rank-one character
is tame. Let $V\to U_1$ be the connected cyclic cover associated
to $\alpha$. Its degree is a power of $p$, hence invertible in
the coefficient field. Inflation--restriction gives an injection
\[
                H^1(U_1,\theta)\longrightarrow H^1(V,\beta|_V).
 \tag{F13.7}
\]

The pullback $\beta|_V$ is tame along $D_Y$. For its unramified
$p$-primary factor this follows by pulling back the cover of
$X_1$ along the proper morphism $Y\to X_1$; for its prime-to-$p$
factor it follows from its prime-to-$p$ image at each DVR.
Every degree-one cocycle for this rank-one system also factors
through the tame quotient. Indeed, the corresponding affine
representation has image in the semidirect product of the additive
coefficient group with the image of $\beta$. On wild inertia
$\beta$ is trivial, so the cocycle restricts to a continuous
homomorphism from a pro-$p$ group to an additive $\ell$-group,
which is zero. The affine cover is therefore tame. Thus (F13.6)
and values of a cocycle on generators give
\[
          \dim H^1(V,\beta|_V)\leq Cd.
 \tag{F13.8}
\]

Now let $M$ be a system occurring in F13. Extend scalars to a
finite field which splits its finite monodromy group. Each simple
constituent $W$ is, by F13a, induced from a rank-one character
$\theta$ on a finite \'etale resolution cover $X_1\to X_0$,
of degree $d=\operatorname{rank}W$. Shapiro's lemma and
(F13.7)--(F13.8) give
\[
 \dim H^1(U_0,W)=\dim H^1(U_1,\theta)
                   \leq C\operatorname{rank}W.
\]
The long exact sequence for a composition series adds these
bounds, proving $\dim H^1(U_0,M)\leq C\operatorname{rank}M$.
Finite extension of the coefficient field preserves these
dimensions; for a fixed finite local system this follows by
tensoring the continuous cochain complexes, or by their filtered
finite-quotient description. Finally apply the fixed-cover
inflation--restriction and complement count already explained
after (F13.5).
\end{proof}

\paragraph{Scope of the reduction.}
The exponent $e$ in F13c is proved uniform. The generator bound
(F13.6) is a separate, unresolved assertion. Bounded degree of a
ramified cover alone has not been shown here to bound the genera
and incidence graph of a resolution linearly in $d$. In particular,
this proposition neither substitutes a bounded-degree statement
for that missing geometric argument nor promotes F13.5 or the
derived-length-three case to Proved.


\subsection{A degree-linear bound for refined-Swan zeros}

The following estimate passed an independent root proof audit.

\begin{proposition}[F13e: zeros on the original exceptional curves]
Retain the fixed pair $X_0,D_0$, matrix $B$, and cutoff vector $t$
of F13. Let $f:X_1\to X_0$ be a connected finite \'etale cover
of degree $d$, and let $\chi$ be a finite $p$-primary character
on $U_1$. Write
\[
 S_\chi=\sum_a s_aC'_a
\]
for its Swan divisor. If $C'_a$ lies above $C_i$ with degree $d_a$
and $s_a>0$, let $Z_a$ be the zero divisor of its nonzero refined
Swan section
\[
 \omega_a\in H^0\!\left(C'_a,
 \Omega^1_{X_1}(\log D_1)|_{C'_a}
                  \otimes\mathcal O_{X_1}(S_\chi)|_{C'_a}\right).
\]
More precisely, $Z_a$ is defined by saturating the image of
$\mathcal O_{C'_a}\xrightarrow{\omega_a}$ in this rank-two
vector bundle. Then
\[
 0\leq\deg Z_a\leq d_a(-B_{ii})(t_i-s_a),\qquad
 \sum_{a:s_a>0}\deg Z_a
       \leq d\sum_i(-B_{ii})t_i.
 \tag{F13.9}
\]
In particular the total number of zeros, counted with their
vanishing multiplicities on the component normalizations, is
bounded linearly in the cover degree. If $s_a=t_i$, that section
has no zero on $C'_a$.
\end{proposition}
\begin{proof}
The section exists and is nonzero by the exact refined-Swan
input used in F10: Krishna--Majumder, Corollary~9.7 and
Theorem~9.4, with the regular $F$-finite hypotheses explicitly
checked there \cite{KrishnaMajumder2026}. Its saturated image
is a line subbundle of the displayed rank-two vector bundle.
The map from $\mathcal O_{C'_a}$ to that line bundle has divisor
$Z_a$, so the degree of the line bundle is $\deg Z_a$.

The residue sequence (F10.4) on $C'_a$ has line subbundle and
line quotient of degrees $2g(C'_a)-2+\nu'_a$ and $0$.
Every line subbundle of its middle term has degree at most
$k'_a=d_a k_i$. Indeed, its map to the line quotient is either
nonzero, in which case its degree is at most zero, or vanishes,
in which case it maps into the line subbundle and has degree at
most $2g(C'_a)-2+\nu'_a$. Tensoring by
$\mathcal O_{X_1}(S_\chi)|_{C'_a}$ therefore gives
\[
                 \deg Z_a\leq d_a k_i+S_\chi\cdot C'_a.
 \tag{F13.10}
\]

The components over a fixed $C_i$ are disjoint: their union is
the finite \'etale inverse image of the smooth curve $C_i$.
The projection formula consequently gives
\[
 (C'_a)^2=d_aB_{ii},\qquad
 \sum_{b:C'_b\mapsto C_j}(C'_b\cdot C'_a)=d_aB_{ij}
                     \quad(j\ne i).
\]
By F13b, $0\leq s_b\leq t_j$ for every component over $C_j$.
All intersection numbers between distinct components are
nonnegative. It follows that
\[
 S_\chi\cdot C'_a
 \leq d_a\left(s_aB_{ii}+\sum_{j\ne i}t_jB_{ij}\right).
\]
The identity $(-B)t=k$ transforms the right side of (F13.10)
into $d_a(-B_{ii})(t_i-s_a)$. This proves the first inequality.
For each $i$, the degrees of all components above $C_i$ sum to
$d$. Dropping the nonnegative $s_a$ terms and summing proves
the second inequality. The final assertion follows because an
effective divisor of degree zero on a proper smooth curve over
an algebraically closed field is zero.
\end{proof}

\paragraph{Exact limitation.}
This estimate concerns the refined-Swan sections on the fixed
resolution and its finite \'etale covers. It does not bound the
number of blowups in a clean-ramification resolution, nor the
cohomology created at those zeros. Those implications require
an additional argument, so (F13.5) and (F13.6) remain open.


\begin{corollary}[F13f: the cutoff below the characteristic]
Suppose the fixed cover and SNC resolution supplied by F10b can
be chosen so that the cutoff vector (F13.1) satisfies $t_i<p$
for every component. Then there is a uniform all-prime
rank-linear $H^1(U,M)$ bound for finite coefficient systems
$M$ with metabelian monodromy. In this subclass, all local
quotients of derived length at most three are uniformly finitely
generated.
\end{corollary}
\begin{proof}
Let $\ell\ne p$. After the fixed cover, each simple constituent
is induced as in F13a from a rank-one character on a finite
\'etale resolution cover. Its $p$-primary part has
Swan conductors bounded by $t_i<p$ under F13b; F13g then forces
all those conductors to be zero. It is unramified on that resolution by
purity. Its prime-to-$p$ part is tame. Induction through the
finite \'etale resolution cover preserves tameness. Thus all
simple constituents are tame. The image of wild inertia is a
finite $p$-group, whose representations in characteristic
$\ell\ne p$ are semisimple, so $M$ itself is tame after the
fixed cover.

Every degree-one cocycle for that pulled-back system factors
through the full tame quotient: its affine monodromy is trivial
on wild inertia, since the original representation is trivial
there and a pro-$p$ group has no nonzero homomorphism to an
additive $\ell$-group. F9 therefore bounds its $H^1$ by
$d_{X_0}\operatorname{rank}M$. Inflation--restriction for the
single fixed cover adds at most
$d(\operatorname{Gal}(U_0/U))\operatorname{rank}M$.

For coefficient $p$, F12 gives the required bound for each
simple constituent; the long exact cohomology sequences for a
composition series give the same rank-linear bound for $M$.
These uniform metabelian coefficient bounds supply precisely
the remaining hypothesis in the complement-count argument
following (F13.5), proving the final assertion.
\end{proof}

This numerical corollary does not assert that a cover satisfying
$t_i<p$ always exists. Nor does it extend the monomial argument
to representations of arbitrary solvable derived length.


\subsection{A residue obstruction and a stronger conductor gap}

The complete argument and exact primary-source formulas passed an
independent root audit.

\begin{lemma}[F13g: a negative Swan intersection forces a $p$-divisible coefficient]
Let $X\to\operatorname{Spec}A$ be an SNC resolution in the
equal-characteristic complete-local surface setting of F10, and
let $S=\sum_i s_iC_i$ be the Swan divisor of a finite
$p$-primary character on the punctured spectrum. If $s_i>0$
and $p\nmid s_i$, then $S\cdot C_i\geq0$. Consequently, if
$S\ne0$, there is a component $C_i$ such that
\[
                 S\cdot C_i<0\quad\hbox{and}\quad p\mid s_i.
 \tag{F13.11}
\]
In particular $S\ne0$ implies $\max_i s_i\geq p$.
\end{lemma}
\begin{proof}
We first check the local residue explicitly. At the completed
DVR of $C_i$, put $n=s_i$ and choose a uniformizer $z$.
The character has a Witt representative
$a=(a_{m-1},\ldots,a_0)$ in the filtration
$p^jv(a_j)\geq-n$. For $j\geq1$, the assumption $p\nmid n$
implies
\[
 p^jv(a_j)\geq-p^j\lfloor n/p^j\rfloor\geq-(n-1).
\]
The refined Swan conductor is represented by
\[
                  -\sum_{j=0}^{m-1}a_j^{p^j-1}\,da_j
 \tag{F13.12}
\]
modulo logarithmic pole order at most $n-1$; these coordinate
conventions and the formula are displayed in
\cite[Definition~2.1 and (2.5), pp.~8--9]{Yatagawa2017Cycle}.
All terms with $j\geq1$ have pole order at most $n-1$.
The remaining coordinate must have pole exactly $n$, since
otherwise the character would have Swan conductor at most
$n-1$. Write $a_0=uz^{-n}$ with $u$ a unit. The logarithmic
residue of $-da_0$ in pole degree $n$ is
$n\bar u$, which is nonzero.

The global refined-Swan section of F10 maps, under the residue
sequence (F10.4) tensored by $\mathcal O_X(S)|_{C_i}$, to a
section of $\mathcal O_X(S)|_{C_i}$. Its generic value is the
nonzero residue just calculated; this generic compatibility is
Krishna--Majumder, Corollary~9.6, and the global map is
Corollary~9.7. Their Lemma~9.11 identifies this local refined
conductor with Kato's refined Swan conductor used in (F13.12)
\cite[Corollaries~9.6--9.7, pp.~55--56, and Lemma~9.11,
p.~57]{KrishnaMajumder2026}.
Thus this line bundle on the proper smooth curve $C_i$ has a
nonzero section. Its degree $S\cdot C_i$ is nonnegative.

Suppose $S\ne0$. Negative definiteness gives
\[
                  0>S^2=\sum_i s_i(S\cdot C_i).
\]
There is therefore an $i$ with $s_i>0$ and $S\cdot C_i<0$.
The first assertion forces $p\mid s_i$, proving (F13.11) and
the final conclusion.
\end{proof}

\begin{corollary}[F13h: improved unramified exponent]
In F13b and F13c, one may choose any integer $e\geq0$ with
\[
                           p^{e+1}>T_0=\max_i t_i.
 \tag{F13.13}
\]
For every rank-one $p$-primary character on every finite \'etale
resolution cover, its $p^e$ multiple then extends unramified
over that cover's resolution. Hence F13c's bounded-order twist
decomposition holds with this improved $e$.
Likewise the $p$-primary inertia kernel in F10 is killed by
$p^e$ whenever $p^{e+1}>\max_i t_i$ for the resolution used
there.
\end{corollary}
\begin{proof}
The Witt restriction step (F10.6) gives every coefficient of the
Swan divisor of $p^e\chi$ at most
$\lfloor T_0/p^e\rfloor<p$. Apply F13g on the resolution of
the complete local normalization underlying the cover. Its
exceptional matrix is negative definite, so that Swan divisor
must be zero. Purity extends the character to the resolution.
The divisible unramified character group from F13c gives the
same twist decomposition with exponent $p^e$.
For F10's abelian inertia kernel, all characters of its
$p$-primary part are killed after multiplying by $p^e$;
character separation gives the last assertion exactly as in
F10.
\end{proof}

\paragraph{Numerical consequence.}
F13f's sufficient hypothesis can therefore be weakened from
$t_i<1$ to $t_i<p$ for every $i$. Under the latter hypothesis,
F13g forces every rank-one inducing character's $p$-primary
part to be unramified; the rest of the proof of F13f is
unchanged. Neither inequality is asserted to hold for an
arbitrary normal surface after the fixed cover.
